\documentclass[journal]{IEEEtran}
\usepackage{amsmath}
\usepackage{graphicx}
\newcommand{\resultfigure}[1]{%
  \IfFileExists{#1}{\includegraphics[width=\columnwidth]{#1}}{%
    \fbox{\parbox[c][0.20\textheight][c]{0.96\columnwidth}{\centering Result figure pending upload.}}%
  }%
}
\usepackage{tikz}
\usetikzlibrary{arrows.meta,positioning,fit,calc,backgrounds}

\ifCLASSINFOpdf

\else

\fi

\hyphenation{op-tical net-works semi-conduc-tor}


\begin{document}

\title{Reinforcement Twinning for Ground-Effect Compensation of a Micro Aerial Vehicle Crossing a Finite Obstacle}

\author{Alessandro Grava and Miguel A. Mendez%
\thanks{A. Grava and M. A. Mendez are with the Department of Control Engineering and System Analysis (SAAS), Universit\'e libre de Bruxelles, Avenue Franklin Roosevelt 50, CP 165/55, 1050 Brussels, Belgium, and also with the von Karman Institute for Fluid Dynamics, Waterloosesteenweg 72, 1640 Sint-Genesius-Rode, Belgium.}%
\thanks{M. A. Mendez is also with the Department of Aerospace Engineering, Universidad Carlos III de Madrid, Avenida de la Universidad 30, 28911 Legan\'es, Madrid, Spain.}}

\markboth{}{}
\maketitle

\begin{abstract}This work presents an implementation of Reinforcement Twinning (RT) applied to a micro aerial vehicle operating in ground effect. The framework couples a Crazyflie~2.1+, the physical environment, with a digital twin (DT) whose ground-effect correction is identified progressively from flight data. During a first information phase, the DT and the remaining uncertainty on the corrective parameters guide the design of information-oriented policies to be deployed on the physical environment. During a subsequent performance phase, model-based optimization first provides a DT-screened policy, while a reward-transfer surrogate combines DT and physical environment rewards to guide data-efficient policy refinement despite residual model error. The resulting policy is implemented as a lightweight spatial reference correction that retains the standard onboard feedback controller. Experiments over a finite obstacle show that the proposed workflow constrains the dominant ground-effect response over the tested clearance range and improves trajectory tracking with a limited number of physical interactions. A proof-of-concept extension to different operating conditions further illustrates how the same workflow can adapt a transferred correction when its performance changes away from the nominal condition. These results support the use of an active DT for model learning and control improvement.\end{abstract}

\begin{IEEEkeywords}Reinforcement Twinning, digital twins, micro aerial vehicles, ground effect, system identification, optimal experiment design, Bayesian optimization, Gaussian processes.\end{IEEEkeywords}

\IEEEpeerreviewmaketitle

\section{Introduction}
\label{Introduction}
\IEEEPARstart{I}{nterest} and investment in unmanned aerial vehicles have grown in recent years. A recent European Commission outlook places the European commercial airborne-drone market at about EUR~14.5 billion by 2030 and indicates that it could exceed EUR~50 billion by 2033~\cite{EuropeanCommission2026}. This growth rests on low operating cost, rapid deployment, and the possibility of replacing or assisting human operators in hazardous missions. Applications range from logistics to emergency response and infrastructure inspection. Small vehicles are particularly useful in cluttered environments that are difficult or dangerous for people to enter. There, micro aerial vehicles (MAVs) with small onboard sensors can map a restricted zone before human operators are sent in.
   
In indoor inspection, monitoring, and emergency response, a MAV can enter confined spaces and cross passages whose clearance is only a few rotor radii. It can approach machinery, beams, pipes, and damaged structures without shutting down a facility or exposing an operator to danger. These missions are also demanding. Satellite positioning is generally unavailable, onboard sensing and computation are limited, and the surrounding structure can degrade communication. Nearby surfaces leave little room for tracking errors and also change the flight dynamics. The propeller wakes interact with these surfaces and alter the forces and moments acting on the vehicle, so the surfaces are part of the aerodynamic environment rather than passive geometric constraints.

Rotor--surface interactions are usually classified as ground, wall, and ceiling effects according to the orientation of the surface relative to the rotor plane~\cite{Matus2020,Carter2021,Franzetti2024}. In a cluttered environment they seldom appear in isolation. A vehicle passing through a narrow tunnel may experience floor, ceiling, and lateral-wall interactions at the same time, with an asymmetric flow field that affects both translation and attitude. Even an isolated MAV rotor responds differently in a wall--ground corner than above a floor~\cite{Prothin2019}. Wall and ceiling effects have received less attention than ground effect, and most available studies concern hovering rotors or vehicles. Wall effect remains weak until the rotor is close to the surface. On a Crazyflie, Carter et al.~\cite{Carter2021} measured a lift change below 5\% near a lateral wall, against up to about 20\% near a floor and 60\% near a ceiling. Ceiling effect is also confined to small clearances, where wake blockage and recirculation increase the force~\cite{Franzetti2024,Cai2024}.

Ground effect (GE), by contrast, has been studied since the 1950s for helicopter rotors. A surface below the rotor redirects the wake radially and reduces the induced velocity and the power required for a given thrust. Cheeseman and Bennett derived the resulting thrust gain of an isolated rotor from its clearance-to-radius ratio and forward speed~\cite{CheesemanBennett1955}, and rotor experiments showed that the near-ground flow changes regime with advance ratio and depends on translational acceleration~\cite{Curtiss1984}. In multirotors, rotor--rotor interaction and the flow beneath the central body change this response~\cite{Matus2020}. It has been described by adding image rotors and an empirical body coefficient to the classical model~\cite{SanchezCuevas2017}, or by exponential, quasi-steady, and Gaussian-process models of the measured thrust~\cite{He2019,HeLeang2020,Franzetti2024ICUAS}. On a tethered Crazyflie~2.0, the near-ground lift increase exceeded twice the classical prediction~\cite{Carter2021}, whereas experiments and computations on other quadrotors report a nonmonotonic thrust or power response near one to one and a half rotor radii~\cite{Tanabe2021,Lei2026}.

Flight near obstacles adds two complications. First, the surface may lie under only some of the rotors. The resulting thrust imbalance produces a moment that a cascaded PID controller rejects poorly~\cite{SanchezCuevas2017}. For a single rotor, the thrust under partial coverage has been correlated with the overlap area between rotor disk and surface~\cite{Yon2026}. Free-flight tests near a building-like edge, however, found that full GE develops only after most of the vehicle has crossed the edge~\cite{Herz2026}. Second, GE is not static. It varies with forward and climb speed~\cite{Kan2019,HeLeang2025} and with the motion of the surface~\cite{He2024}, and dynamic quadrotor tests have complemented the static characterizations~\cite{DelContBernard2017}.

Control studies compensate GE online or learn it from flight data. Adaptive control~\cite{Wei2019} and disturbance observers~\cite{He2019,LiX2024} counter the near-ground force, and Li et al.~\cite{LiX2024} also fitted a GE model to observer estimates recorded in hovering flights. Deep networks trained on flight data~\cite{Shi2019} and richer models that include surface motion and past states~\cite{He2024} or GE torques and rotor drag~\cite{Yang2025} provide feedforward or model-based compensation. In aircraft, GE coefficients have been identified from flight data recorded without dedicated maneuvers~\cite{Vitale2010}. Flight-vehicle system identification, by contrast, designs the inputs to maximize the information on the unknown parameters~\cite{Mehra1974,Jategaonkar2015,Morelli2016}. Robot trajectories have been synthesized to maximize the Fisher information~\cite{Wilson2014}, and informative trajectories have been planned for other multirotor parameters~\cite{Bahnemann2017}. To our knowledge, no study has selected the flights of a multirotor by their information on the GE parameters.

Accurate GE prediction matters in cluttered flight because an aerodynamic model must do more than reproduce a measured force. It must let the controller keep the vehicle on a collision-free path from one obstacle to the next. A thrust change above one surface can produce a vertical excursion that is harmless locally but uses up the clearance available at a following beam, ceiling, or doorway. The interaction must therefore be predicted and compensated for before a tracking error develops. Finite obstacles make this requirement explicit. The vehicle enters GE, remains in it briefly, and then crosses a trailing edge, so leading-edge exposure, transient wake development, and post-obstacle recovery all occur within one mission.

Most MAVs rely on cascaded proportional-integral-derivative (PID) loops. They are cheap to compute and robust near their nominal operating point, which suits small onboard processors. However, they are reactive. An unmodeled disturbance must first produce an error before feedback counters it, and close to an obstacle this delay can consume the available safety margin. Predictive methods instead modify the command or reference before the disturbance develops. Model predictive control optimizes future behavior while accounting for dynamics and constraints~\cite{MPCReview2021}, also on quadrotors with learned residual dynamics~\cite{Torrente2021}. Its effectiveness depends on the prediction model and on the cost of solving the optimization online. An incomplete model can produce a wrong correction, whereas a more faithful aerodynamic model may be difficult to identify and too expensive to optimize repeatedly on the limited hardware of a MAV.

Reinforcement Twinning (RT) was introduced to address this trade-off. It couples a physical system, called the physical environment, with an active physics-based digital twin (DT)~\cite{SchenaRT2023,Mendez2026,PolettiRT2025}. A conventional DT is often understood as a virtual representation on which the behavior of a physical environment can be reproduced. An active DT additionally learns from physical observations and uses the updated representation to influence subsequent operation. It acts as both a \emph{steering representation} and an \emph{acting model} in the decision process~\cite{Korenhof2021,Wagg2025}. Experimental data refine its uncertain physics, and virtual rollouts select the next policy deployed in the physical environment. RT has been applied to wind-turbine control, gust rejection on a flapping-wing vehicle, and cryogenic thermal management~\cite{SchenaRT2023}, to the hybrid control of a flapping-wing drone~\cite{PolettiRT2025}, and to a flow-control benchmark~\cite{Mendez2026}, all in simulation.

Early RT algorithms combined model-based and model-free control. A model-based policy supplied an efficient, physically informed starting point, and a model-free policy kept improving when the model-based one reached a local or structurally biased optimum. A referee switched between the two according to their observed or predicted performance. When learning stagnated, policy cloning copied the weights of the better policy to the other learner, so that the search resumed from a better region~\cite{SchenaRT2023,PolettiRT2025}. The model-based branch thus made the first interactions less risky and accelerated early learning, without tying the final controller to the model-based optimum.

More recent RT formulations cast the exchange between model-based and data-driven learning as a twin-calibrated Bayesian policy search~\cite{Mendez2026}, in which each physical episode has one of two purposes. An information-oriented policy is chosen to improve the uncertain physics in the DT, and a performance-oriented policy to improve control. The two objectives are coupled through the accumulated data and the current DT, but they need not be merged into one scalar objective. This distinguishes RT from dual control, which balances the regulation and information consequences of each action~\cite{Feldbaum1960,Filatov2000,Mesbah2018}. RT instead assigns the two roles to separate episodes, linked by the DT, which is updated by replaying the recorded episodes (factual replay). Once the physical data constrain the DT well enough for the task, the DT solves the model-based control problem offline and screens candidate policies against state and clearance constraints before deployment. The resulting model-based policy accounts for the dynamics learned from the experiments. If it is expressed as a compact feedforward reference correction, the standard feedback controller is left unchanged and still stabilizes the vehicle.

Identification reduces parametric uncertainty but cannot remove the structural discrepancy between a compact physics model and the physical environment. Additional states and correction terms may reduce this model discrepancy, but a DT that remains cheap to optimize still omits part of the physical response. RT therefore also learns how control performance transfers from the DT to the physical environment. Each tested policy yields a DT reward and a physical reward. In the resulting reward-transfer surrogate, an affine relation captures the global calibration between the two rewards, and a Gaussian process (GP) represents the remaining discrepancy, whose uncertainty is largest in sparsely explored regions. This is the autoregressive multi-fidelity structure of Kennedy and O'Hagan~\cite{Kennedy2000,Kennedy2001,Forrester2007,Rasmussen2006}.

The surrogate turns the last stage of RT into a Bayesian policy search. Expected Improvement (EI) weighs the predicted physical reward against the surrogate uncertainty and proposes a candidate that is promising, insufficiently explored, or both~\cite{Mockus1975,Jones1998}. The physical reward of the tested candidate updates the surrogate, and the search is repeated. The search can therefore move beyond the model-based policy even if the DT retains a systematic bias. Unlike a model-free search, however, it starts from the model-based policy and stays within a neighborhood that the DT has already simulated and screened. The DT supplies the initial direction and the admissible set, and physical data correct its ranking where needed.

The repeated mission also invites comparison with two established approaches. Iterative learning control refines a feedforward signal from the tracking errors of previous trials of the same task~\cite{Bristow2006}, and it has been applied to quadrotor trajectory tracking~\cite{Schoellig2012,Meraglia2021}. Bayesian optimization has tuned quadrotor controllers directly on hardware~\cite{Berkenkamp2016} and has combined simulations and physical experiments as information sources of different cost~\cite{Marco2017}. RT differs from both in that the flights first identify an interpretable twin, which then designs and screens the policies, and the physical data correct the twin's ranking of candidate policies.

This paper presents an experimental demonstration of a two-stage form of RT on a physical MAV. A Crazyflie~2.1+ crosses a finite horizontal obstacle at constant forward speed, so that GE appears and disappears within one mission. The mission is repeated over an obstacle of known position and geometry, with external positioning, which isolates the GE transient from perception and planning. The initial DT is the six-degree-of-freedom Crazyflie model of Richardson and Pretorius~\cite{Richardson2025}, with an added static GE law whose two parameters start from the Cheeseman--Bennett values. This low-order structure can be identified and used for control design, but it leaves a model discrepancy, mainly at the obstacle edges.

In the first stage, E-optimal information-oriented flights identify the two GE parameters, which are then frozen. In the second stage, the DT computes a model-based policy, applied as a spatial reference correction above the unchanged onboard PID controller. Physical rewards of nearby policies calibrate the reward-transfer surrogate, and EI searches for further gains. The two stages are consecutive rather than interleaved, so that the role of each flight is explicit. Finally, the correction learned at the nominal condition is transferred to two other combinations of clearance and speed, and a short EI search is tested at one of them.

The paper makes three contributions. First, it applies RT to a physical vehicle. A differentiable DT of a Crazyflie~2.1+ and its onboard controller designs both the identification flights and the control reference. Second, E-optimal flights over two clearance domains identify a static GE law whose amplitude is about twice the Cheeseman--Bennett value, and the DT yields a reference correction that raises the physical tracking reward from $0.544$ to $0.855$ without changing the onboard controller. Third, after nine EI-selected policies, a reward-transfer surrogate finds a correction with physical reward $0.952$ that the DT ranks as nearly equivalent to the model-based one. The transfer of this correction to two off-nominal conditions is also tested.

The remainder of the paper is organized as follows. Section~\ref{sec:methodology} describes the RT loop, the physical environment and DT, and the test campaign; Section~\ref{sec:results} reports the experiments; and Section~\ref{sec:conclusion} states the conclusions and limitations.

\section{Methodology}\label{sec:methodology}


\subsection{Reinforcement Twinning Framework}
\label{sec:rt-framework}

In RT, the physical environment and an active, physics-based DT evolve side by side with complementary objectives~\cite{SchenaRT2023,Mendez2026}. Physical data update the DT and score the policies it designs, and the DT proposes the next experiments and policies.

\paragraph*{The two environments and their objective}
The physical environment is denoted by a filled circle $(\bullet)$ and the DT by an empty circle $(\circ)$, following the black- and white-box symbology, respectively. A physical interaction is indexed by $n$, while $k$ denotes the time step within that interaction. The hidden physical state $\mathbf s_{\bullet,k}^{(n)}$ evolves according to
\begin{equation}
\mathbf s_{\bullet,k+1}^{(n)}=
\mathbf f_{\bullet}\!\left(
\mathbf s_{\bullet,k}^{(n)},
\mathbf u_{\bullet,k}^{(n)},
\mathbf d_k^{(n)}
\right),
\label{eq:rt-physical-state}
\end{equation}
where $\mathbf u_{\bullet,k}^{(n)}$ is the command generated under the deployed policy and $\mathbf d_k^{(n)}$ collects exogenous variables and disturbances that are not chosen by the policy. These may include environmental loads, variations of the available actuation, imperfect initial conditions, or disturbances in the measurement chain. The state is not assumed to be directly available as in a black-box framework. What enters the RT loop is the noisy observation
\begin{equation}
\mathbf o_{\bullet,k}^{(n)}=
\mathbf h_{\bullet}\!\left(\mathbf s_{\bullet,k}^{(n)}\right)
+\boldsymbol{\nu}_{k}^{(n)},
\label{eq:rt-physical-observation}
\end{equation}
with $\boldsymbol{\nu}_{k}^{(n)}$ representing measurement noise and unmodeled observation errors. The complete experimental record also contains the deployed policy, the applied inputs, the relevant operating conditions, and the physical reward $\mathcal R_{\bullet}^{(n)}$.

The DT reproduces the same state evolution with a parameterized transition model,
\begin{equation}
\mathbf s_{\circ,k+1}^{(n)}=
\mathbf f_{\circ}\!\left(
\mathbf s_{\circ,k}^{(n)},
\mathbf u_{\circ,k}^{(n)},
\mathbf d_k^{(n)};
\boldsymbol{\theta}_{p}^{(m)}
\right),
\label{eq:rt-digital-state}
\end{equation}
and generates the corresponding DT observation
\begin{equation}
\mathbf o_{\circ,k}^{(n)}=
\mathbf h_{\circ}\!\left(
\mathbf s_{\circ,k}^{(n)};
\boldsymbol{\theta}_{p}^{(m)}
\right),
\label{eq:rt-digital-observation}
\end{equation}
where the vector $\boldsymbol{\theta}_{p}^{(m)}$ contains the unknown or uncertain model parameters to be identified and $m$ counts their updates. Every physical episode increments $n$, whereas $m$ advances only when model identification is active; therefore, $m \le n$. During factual replay, the DT receives the recorded policy, exogenous inputs, operating condition, and a reconstructed initial state. The mismatch between $\mathbf o_{\circ}$ and $\mathbf o_{\bullet}$
then measures how much of the physical response is not explained by the current twin. During virtual evaluation, the same DT predicts the outcome of policies that have not yet been tested physically.

Each environment has its own objective. The physical environment seeks policy parameters $\boldsymbol{\theta}_{\pi}$ in the domain $\Theta_{\pi}$ that maximize the expected physical reward,
\begin{equation}
\boldsymbol{\theta}_{\pi}^{\star}
=\arg\max_{\boldsymbol{\theta}_{\pi}\in\Theta_{\pi}}
\mathrm{E}\!\left[\mathcal R_{\bullet}
(\boldsymbol{\theta}_{\pi})\right],
\label{eq:rt-physical-objective}
\end{equation}
whereas the digital twin searches for model parameters $\boldsymbol{\theta}_{p}$  in the domain $\Theta_{p}$ that make factual replay consistent with the accumulated observations,
\begin{equation}
\widehat{\boldsymbol{\theta}}_{p}^{(m+1)}
=\arg\min_{\boldsymbol{\theta}_{p}\in\Theta_p}
\mathcal L_p\!\left(
\mathcal B_n;\boldsymbol{\theta}_{p}
\right),
\label{eq:rt-model-objective}
\end{equation}
where $\mathcal B_n=\{\mathcal E^{(1)},\ldots,\mathcal E^{(n)}\}$ is the replay buffer of physical episode records
$
\mathcal{E}^{(n)} = \left\{
\boldsymbol{o}_{\bullet}^{(n)},
\boldsymbol{\theta}_{\pi}^{(n)},
e_n,
\mathbf d^{(n)},
\mathcal{R}_{\bullet}^{(n)}
\right\}$,
and $\mathcal{L}_p$ compares the DT and physical observations selected for identification.

\paragraph*{Policies and interaction schedule}
The interaction between the two environments is mediated by control policies that are designed within the DT and subsequently deployed in the physical environment. The resulting physical episode then provides the experimental evidence required by the learning process. Depending on the current objective, the DT designs either an information-oriented or a performance-oriented policy. The information-oriented policy $\pi_{\iota}(\cdot\mid\boldsymbol{\theta}_{\iota})$, with $\boldsymbol{\theta}_{\iota}\in\Theta_{\iota}$, is designed to generate observations that are informative about the model parameters $\boldsymbol{\theta}_{p}$. Conversely, the performance-oriented policy $\pi_{\bullet}(\cdot\mid\boldsymbol{\theta}_{\pi})$, with $\boldsymbol{\theta}_{\pi}\in\Theta_{\pi}$, is designed using the current DT to maximize the expected physical reward. The two families are defined by the objective used to design them and may have different parameterizations. However, the data generated by each physical episode can still support both learning objectives. A performance-oriented policy may excite dynamics that improve the model-parameter estimate, while an information-oriented policy also yields a physical reward. The present campaign does not use these rewards, because the information-oriented references depart from the nominal clearance.
A binary selector $e_n\in\{0,1\}$ determines which policy is designed and deployed in physical episode $n$,
\begin{equation}
\pi_{\bullet}^{(n)}=
\begin{cases}
\pi_{\iota}\!\left(\cdot\mid\boldsymbol{\theta}_{\iota}^{(n)}\right),
& e_n=1,\quad \text{information-oriented},\\[2mm]
\pi_{\bullet}\!\left(\cdot\mid\boldsymbol{\theta}_{\pi}^{(n)}\right),
& e_n=0,\quad \text{performance-oriented}.
\end{cases}
\label{eq:rt-policy-schedule}\end{equation}

In the general RT formulation, these two types of interaction are freely alternated during the campaign as the twin and policy evolve. For simplicity, the present work performs all information-oriented episodes first, until the model meets the stopping criterion of the information phase. The model parameters are then frozen, and the campaign proceeds with performance-oriented episodes. This defines an \emph{information phase} followed by a separate \emph{performance phase}.

\paragraph*{Information phase}
In the information phase, $\boldsymbol{\theta}_{p}$ is iteratively updated using the accumulated experimental data, after which the updated DT designs the next information-oriented policy. Given the current replay buffer, Eq.~\eqref{eq:rt-model-objective} is solved until the model update has converged. The updated DT then generates a set of admissible candidate interactions within the information-oriented policy domain. The DT predictions along these interactions are differentiated with respect to $\boldsymbol{\theta}_p$. As $\mathbf S(\boldsymbol{\theta}_{\iota})=\partial\mathbf o_{\circ}/\partial\boldsymbol{\theta}_p$ denotes the sensitivity of the predicted observations and $\boldsymbol{\Sigma}_{\nu}$ their noise covariance, the information expected after a candidate policy is
\begin{equation}
\mathbf F_{n+1}(\boldsymbol{\theta}_{\iota})
=
\mathbf F_n+
\mathbf S^{\top}(\boldsymbol{\theta}_{\iota})
\boldsymbol{\Sigma}_{\nu}^{-1}
\mathbf S(\boldsymbol{\theta}_{\iota}),
\label{eq:rt-fisher-update}
\end{equation}
where $\mathbf F_n$ is the Fisher information accumulated from the experiments already stored in $\mathcal B_n$. The policy deployed in episode $n+1$ is selected with an E-optimal criterion~\cite{Mehra1974,Goodwin1977,Franceschini2008},
\begin{equation}
\boldsymbol{\theta}_{\iota}^{(n+1)}
=
\arg\max_{\boldsymbol{\theta}_{\iota}\in\Theta_{\iota,\mathrm{adm}}}
\lambda_{\min}\!\left[
\mathbf F_{n+1}(\boldsymbol{\theta}_{\iota})
\right].
\label{eq:rt-e-optimal}
\end{equation}
Maximizing the smallest eigenvalue targets the least-observed joint parameter direction, so the E-optimal criterion already favors novelty in the information collected. The admissible set $\Theta_{\iota,\mathrm{adm}}$ is obtained through virtual rollouts and may enforce state, input, rate, and clearance constraints.

The selected policy is deployed in the physical environment, its record is appended to $\mathcal B_n$, and the model is identified again from the enlarged dataset. Two convergence levels are distinguished. Each parameter update is complete only when the identification loss, the model parameters, and the projected gradient remain within prescribed tolerances over a monitoring window. The information phase ends when new episodes no longer change the converged estimates appreciably or reduce the magnitude of their correlation. In the present campaign this was judged from the parameter history, the correlation, and the replay of the baseline episode (Section~\ref{sec:results-inf}); no numerical threshold was set in advance.

The converged model parameters, denoted by $\widehat{\boldsymbol{\theta}}_p$, are then used to compute the model-based policy $\pi_{\mathrm{MB}}$, subject to the reference-amplitude and clearance constraints checked in the DT,
\begin{equation}
\boldsymbol{\theta}_{\pi,\mathrm{MB}}
=
\arg\max_{\boldsymbol{\theta}_{\pi}\in\Theta_{\pi,\mathrm{adm}}}
\mathcal R_{\circ}\!\left(
\boldsymbol{\theta}_{\pi};
\widehat{\boldsymbol{\theta}}_p
\right).
\label{eq:rt-mb-optimum}
\end{equation}
In the present work, $\pi_{\mathrm{MB}}$ marks the transition between the two phases. In a more general RT approach, the model-based policy can be updated after every $\boldsymbol{\theta}_p$ update.

\paragraph*{Performance phase}
In the present work, the performance phase starts from $\boldsymbol{\theta}_{\pi,\mathrm{MB}}$ and keeps $\widehat{\boldsymbol{\theta}}_p$ frozen. The search is restricted to a neighborhood of $\pi_{\mathrm{MB}}$ because this model-based policy is expected to improve upon the baseline policy and provides a physically informed reference. The objective is therefore to improve upon $\pi_{\mathrm{MB}}$ by learning the difference between the DT reward and the physical reward. A small set of policies around $\pi_{\mathrm{MB}}$ is first screened in the DT and then flown. These paired rewards initialize the reward-transfer surrogate,
\begin{equation}
\begin{aligned}
\mathcal R_{\bullet}(\boldsymbol\xi)
&=\gamma_0+
\gamma_1\mathcal R_{\circ}
(\boldsymbol\xi;\widehat{\boldsymbol{\theta}}_p)
+\delta(\boldsymbol\xi)+\varepsilon,\\
\delta&\sim\mathcal{GP}(0,k),\qquad
\varepsilon\sim\mathcal N(0,\sigma_{\varepsilon}^{2}).
\end{aligned}
\label{eq:rt-surrogate}
\end{equation}
The descriptor $\boldsymbol\xi=[\boldsymbol{\theta}_{\pi}^{\top},\boldsymbol\chi^{\top}]^{\top}$ contains the policy parameters and the operating-condition vector $\boldsymbol\chi$ that describes the performance task. The intercept $\gamma_0$ and slope $\gamma_1$ describe the global affine relation between the DT and physical rewards. The Gaussian-process (GP) discrepancy $\delta(\boldsymbol\xi)$ models the policy-dependent part of the transfer that the affine relation does not explain. Its Gaussian-process prior assigns similar discrepancies to nearby descriptors through the kernel $k(\boldsymbol\xi,\boldsymbol\xi')$, while $\varepsilon$ represents experimental and reward-evaluation noise. The present campaign reduces $\sigma_{\varepsilon}$ to a numerical jitter (Section~\ref{sec:performance-phase}).
The GP is conditioned on the residuals between the physical rewards $\mathcal R_{\bullet}$ and the corresponding affine predictions obtained from the DT rewards $\mathcal R_{\circ}$ collected in the replay buffer $\mathcal B_n$. Its posterior provides the mean GP discrepancy $\mu_{\delta}(\boldsymbol\xi)$ and its standard deviation $\sigma_{\delta}(\boldsymbol\xi)$. The predicted physical reward therefore has mean
\begin{equation}
\mu_{\bullet}(\boldsymbol\xi)=
\gamma_0+
\gamma_1\mathcal R_{\circ}
(\boldsymbol\xi;\widehat{\boldsymbol{\theta}}_p)
+\mu_{\delta}(\boldsymbol\xi)
\label{eq:rt-surrogate-mean}
\end{equation}
and GP interpolation uncertainty $\sigma_{\bullet}(\boldsymbol\xi)=\sigma_{\delta}(\boldsymbol\xi)$.

The next performance-oriented policy is selected through Expected Improvement. Let $\mathcal R_{\bullet}^{+}$ denote the best physical reward observed so far and $\zeta\geq0$ an optional exploration margin. For $\sigma_{\bullet}(\boldsymbol\xi)>0$, define $Z(\boldsymbol\xi)=[\mu_{\bullet}(\boldsymbol\xi)-\mathcal R_{\bullet}^{+}-\zeta]/\sigma_{\bullet}(\boldsymbol\xi)$, so that
\begin{equation}
\operatorname{EI}(\boldsymbol\xi)=
\left[\mu_{\bullet}(\boldsymbol\xi)-\mathcal R_{\bullet}^{+}-\zeta\right]
\Phi\!\left[Z(\boldsymbol\xi)\right]
+
\sigma_{\bullet}(\boldsymbol\xi)
\phi\!\left[Z(\boldsymbol\xi)\right],
\label{eq:rt-expected-improvement}
\end{equation}
where $\Phi$ and $\phi$ are the standard-normal cumulative distribution and density functions. EI is set to zero when $\sigma_{\bullet}(\boldsymbol\xi)=0$. Its first term favors policies expected to exceed the best physical reward, while the second favors uncertain candidates. The next interaction is
\begin{equation}
\boldsymbol\xi^{(n+1)}
=
\arg\max_{\boldsymbol\xi\in\mathcal X_{\mathrm{adm}}}
\operatorname{EI}(\boldsymbol\xi).
\label{eq:rt-next-performance}
\end{equation}
After deployment, the physical reward is added to $\mathcal B_n$ and the reward-transfer surrogate is updated.

\paragraph*{Summary of the Reinforcement Twinning loop}
Fig.~\ref{fig:rt_framework} shows one iteration of the RT loop, from the deployed policy to the choice of the next one. The selector $e_n$ activates one of two learning objectives without changing the architecture of the loop. An information-oriented episode modifies the physical model through factual replay, whereas a performance-oriented episode leaves that model frozen and updates the reward-transfer surrogate.

\begin{figure*}[t]
\centering
\resizebox{0.98\textwidth}{!}{%
\begin{tikzpicture}[
  >=Latex,
  font=\small,
  block/.style={draw=black!75, rounded corners=2pt, align=center,
                text width=6.15cm, minimum height=1.55cm,
                inner sep=5pt, fill=white},
  flow/.style={->, thick, draw=black!85},
  note/.style={font=\footnotesize, align=center, inner sep=1.5pt},
  number/.style={circle, draw=black, fill=white, inner sep=1.5pt,
                 minimum size=5.5mm, font=\small}
]

% Panel headers and images
\node[font=\bfseries\normalsize,align=center] (ptitle) at (0,7.70)
  {Physical environment ($\bullet$)\\[-0.2mm]{\normalfont\normalsize Objective: improve policy parameters $\theta_\pi^{(n)}$}};
\node (pimage) at (0,5.55)
  {\includegraphics[height=3.08cm]{Methodology/Physical_Environment.png}};

\node[font=\bfseries\normalsize,align=center] (dtitle) at (9.0,7.70)
  {Digital twin ($\circ$)\\[-0.2mm]{\normalfont\normalsize Objective: improve model parameters $\theta_p^{(m)}$}};
\node (dimage) at (9.0,5.55)
  {\includegraphics[height=3.08cm]{Methodology/Digital_Twin.png}};

% Physical environment
\node[block] (b1) at (0,2.75) {
  \textbf{Deployed policy $\pi_\bullet^{(n)}$}\\[1mm]
  $e_n=1$: information $\pi_\iota(\cdot\mid\theta_\iota^{(n)})$\\
  $e_n=0$: performance $\pi_\bullet(\cdot\mid\theta_\pi^{(n)})$
};

\node[block] (b2) at (0,0.35) {
  \textbf{Physical episode $n$}\\[1mm]
  hidden state $\mathbf{s}_\bullet$, input $\mathbf d$\\
  noisy observations $\mathbf{o}_\bullet$
};

\node[block] (b3) at (0,-2.10) {
  \textbf{Record $\mathcal E^{(n)}$}\\[1mm]
  $\mathcal E^{(n)}$ put in replay buffer $\mathcal B$;\\
  evaluated $\mathcal R_\bullet^{(n)}$
  
};

% Digital twin
\node[block] (b6) at (9.0,2.75) {
  \textbf{Choice of the next interaction}\\[1mm]
  $e_n=1$: $\mathcal J_{\mathrm{info}}\rightarrow\theta_\iota^{(n+1)}$
  \\
  $e_n=0$: $\widehat{\mathcal R}_\bullet(\mathcal R_\bullet,\mathcal R_\circ,\theta_\pi)$
  $\rightarrow\theta_\pi^{(n+1)}$
};

\node[block] (b5) at (9.0,0.35) {
  \textbf{Virtual candidate evaluation}\\[1mm]
  $e_n=1$: informativeness $\mathcal J_{\mathrm{info}}(\theta_\iota)$
  \\
  $e_n=0$: $\theta_p^{(m+1)}\rightarrow\theta_{\pi,\mathrm{MB}}$
};

\node[block] (b4) at (9.0,-2.10) {
  \textbf{Factual replay}\\[1mm]
  
  initial-state reconstruction\\
  $e_n=1$: $\mathbf o_\circ$ vs. $\mathbf o_\bullet\rightarrow\theta_p^{(m+1)}$ updated\\
  $e_n=0$: $\theta_p^{(m+1)}=\theta_p^{(m)}$
};

% Number labels
\node[number, anchor=center] at ([xshift=-3.08cm,yshift=0.78cm]b1.center) {1};
\node[number, anchor=center] at ([xshift=-3.08cm,yshift=0.78cm]b2.center) {2};
\node[number, anchor=center] at ([xshift=-3.08cm,yshift=0.78cm]b3.center) {3};
\node[number, anchor=center] at ([xshift=-3.08cm,yshift=0.78cm]b4.center) {4};
\node[number, anchor=center] at ([xshift=-3.08cm,yshift=0.78cm]b5.center) {5};
\node[number, anchor=center] at ([xshift=-3.08cm,yshift=0.78cm]b6.center) {6};

% Main loop
\draw[flow] (b1.south) -- node[note,right]{deploy} (b2.north);
\draw[flow] (b2.south) -- node[note,right]{measure} (b3.north);
\draw[flow] (b3.east) -- node[note,below]{replay buffer $\mathcal B$} (b4.west);
\draw[flow] (b4.north) -- node[note,right,text width=3.35cm,align=center]{
  $\theta_p^{(m+1)}$\\
  } (b5.south);
\draw[flow] (b5.north) -- node[note,right]{
  $\mathcal R_\circ,\,\theta_{\pi,\mathrm{MB}},\,\mathcal J_{\mathrm{info}}$} (b6.south);
\draw[flow] (b6.west) -- node[note,above,text width=3.2cm,align=center]{
  next interaction:\\[0.8mm] $\theta_\pi^{(n+1)}$ or $\theta_\iota^{(n+1)}$} (b1.east);

% Measured performance data bypasses factual identification
\draw[flow] (b3.east) -- ++(1.00,0) coordinate (retcorner)  -- node[note,rotate=90,above]{
  measured returns $(\theta_\pi^{(n)},\mathcal R_\bullet^{(n)})$} (retcorner |- b6.south west)  -- (b6.south west);

% Panel backgrounds
\begin{scope}[on background layer]
  \node[fit=(ptitle)(pimage)(b1)(b2)(b3),
        fill=black!5, draw=black!12, rounded corners=8pt,
        inner sep=10pt] {};
  \node[fit=(dtitle)(dimage)(b4)(b5)(b6),
        fill=white, draw=black!35, rounded corners=8pt,
        inner sep=10pt] {};
\end{scope}

\end{tikzpicture}%
}
\caption{RT loop adopted in this work. The physical environment and its active DT evolve side by side through physical observations and virtual candidate evaluations. The selector $e_n$ distinguishes information-oriented ($e_n=1$) and performance-oriented ($e_n=0$) episodes within a common loop. In the figure, $\mathcal J_{\mathrm{info}}$ is the E-optimal criterion of Eq.~\eqref{eq:rt-e-optimal} and $\widehat{\mathcal R}_\bullet$ the surrogate prediction of the physical reward, written $\mu_\bullet$ in Eq.~\eqref{eq:rt-surrogate-mean}.}
\label{fig:rt_framework}
\end{figure*}
The blocks of Fig.~\ref{fig:rt_framework} perform the following operations.
\begin{enumerate}
    \item Policy deployment. The policy selected in the previous iteration is deployed; the first episode uses the baseline policy $\pi_{\mathrm{BL}}$, which sets the reference performance.
    \item Physical episode. The physical environment evolves under its hidden dynamics, the deployed policy, and the exogenous input $\mathbf d$, and returns the noisy observations $\mathbf o_{\bullet}$ of Eq.~\eqref{eq:rt-physical-observation}.
    \item Recording. The record $\mathcal E^{(n)}$ is appended to the replay buffer $\mathcal B$, so that every model update uses the complete dataset.
    \item Factual replay. For $e_n=1$, the DT replays the recorded episode, and the mismatch between $\mathbf o_{\circ}$ and $\mathbf o_{\bullet}$ updates $\boldsymbol{\theta}_p^{(m)}$ to $\boldsymbol{\theta}_p^{(m+1)}$. For $e_n=0$, the model stays frozen, and the new policy--reward pair enters the reward-transfer surrogate.
    \item Virtual candidate evaluation. The DT simulates candidate policies. Information-oriented candidates are scored by their predicted information on $\boldsymbol{\theta}_p$. Performance-oriented candidates around $\pi_{\mathrm{MB}}$ that pass the DT checks receive a DT reward and a surrogate prediction of their physical reward. Many candidates are evaluated virtually, but only the selected one is flown.
    \item Choice of the next interaction. Equation~\eqref{eq:rt-e-optimal} ($e_n=1$) or Eq.~\eqref{eq:rt-next-performance} ($e_n=0$) selects the next policy, $n$ advances to $n+1$, and the loop returns to Block~1.
\end{enumerate}
\subsection{Test Case Physical and Digital Environments}
\label{sec:setup}
\paragraph*{Experimental setup}
The experiments are designed to reproduce a transition between out-of-ground-effect (OGE) and in-ground-effect (IGE) conditions over a finite obstacle while maintaining a constant forward flight speed. The experimental setup and adopted reference frame are illustrated in Fig.~\ref{fig:experimental_setup}.

The campaign is performed using a Crazyflie 2.1+ MAV from Bitcraze flying in the $x$-$z$ plane at $y=0$~m. The horizontal axes keep the origin $x=y=0$ of the positioning frame, whereas the vertical origin is the upper surface of the obstacle, a table whose top is $0.63$~m above the floor. Positive $z$ is the height above this surface of the rotor plane, which lies $8.5$~mm above the center of mass and $1.5$~mm above the Lighthouse deck. With $R=23$~mm the rotor radius, $z/R$ is the nondimensional clearance. After takeoff, the MAV moves to $x=0.20$~m, climbs to the desired clearance $z^\ast$, and hovers there for 10~s. It then accelerates smoothly and reaches the desired forward speed around $x=1.0$--$1.2$~m while maintaining the prescribed $z/R$. The region of interest used for comparison with the DT results spans $x=1.70$--$4.00$~m. The MAV subsequently decelerates and lands around $x=4.20$~m. Within the region of interest, the obstacle extends from $x=2.20$~m to $x=3.01$~m, a streamwise length of $0.81$~m, and it is $1.0$~m wide. In the results, the streamwise coordinate is shifted to the beginning of the region of interest and normalized by $R$; consequently, this region covers $0\leq x/R\leq100$, while the obstacle occupies approximately $21.7\leq x/R\leq57.0$.

The Crazyflie~2.1+ carries an onboard inertial measurement unit (IMU) for attitude control and a Lighthouse positioning deck for position control. With its 350~mAh battery, its measured mass is $m=32.42$~g. For continuous positioning, the Lighthouse deck receives the infrared sweeps of two Lighthouse base stations. Their placement is selected such that the MAV remains visible to at least one station throughout the complete trajectory. In the region directly above the obstacle, which constitutes the main area of interest for the GE experiments, the fields of view of the two stations intentionally overlap. This overlap reduces the risk of losing the position estimate when one optical link is interrupted. For Lighthouse~V2, Taffanel et al.~\cite{Taffanel2021Lighthouse} reported mean Euclidean errors of 2--4~cm in flight, with outliers up to about 5~cm, and a stationary jitter of 0.3--0.7~mm depending on the estimator. 

\begin{figure*}[t]
    \centering
    \resizebox{0.98\textwidth}{!}{%
    \begin{tikzpicture}[>=Latex,font=\small]
        \path[use as bounding box] (-7.4,-2.15) rectangle (7.4,5.25);

        % Lighthouse fields of view
        \begin{scope}[on background layer]
            \clip (-7.4,-2.15) rectangle (7.4,5.25);
            \fill[red!18,opacity=0.16] (-6.25,3.95) -- (-8.00,-2.05) -- (23.00,-2.05) -- cycle;
            \fill[red!18,opacity=0.16] ( 6.25,3.95) -- ( 8.00,-2.05) -- (-23.00,-2.05) -- cycle;
        \end{scope}

        % Lighthouse stations; station 1 is mirrored from the same PNG
        \node (lh1) at (-6.25,4.15) {\reflectbox{\includegraphics[width=1.45cm]{Setups/lighthouse.png}}};
        \node (lh2) at ( 6.25,4.15) {\includegraphics[width=1.45cm]{Setups/lighthouse.png}};
        \node[draw,rounded corners=2pt,fill=white,inner sep=3pt,anchor=south] at (-6.25,4.86) {Lighthouse station 1};
        \node[draw,rounded corners=2pt,fill=white,inner sep=3pt,anchor=south] at ( 6.25,4.86) {Lighthouse station 2};

        % Nominal crossing and repeated vehicle poses
        \draw[densely dashed,black!55,line width=0.8pt] (-5.45,1.55) -- (5.45,1.55);
        \node[green!45!black,font=\normalsize] at (-3.25,3.18) {$v_x^\ast/v_h\simeq\mathrm{const}$};
        \draw[->,very thick,green!45!black] (-3.90,2.88) -- (-2.60,2.88);
        \foreach \xx/\op in {-4.65/0.18,-2.45/0.34,-0.20/0.52,2.05/0.74,4.30/1.00}{
            \node[opacity=\op] at (\xx,1.62) {\includegraphics[width=2.20cm]{Setups/drone.png}};
        }
        \node[draw,rounded corners=2pt,fill=white,inner sep=3pt,align=center,anchor=south] at (0.85,2.38)
            {Crazyflie 2.1+ with Lighthouse deck};

        % Finite obstacle and vertical clearance
        \node[anchor=north] (obs) at (0,0.62) {\includegraphics[width=6.60cm]{Setups/obstacle.png}};
        \node[draw,rounded corners=2pt,fill=white,inner sep=3pt] at (0,-1.45) {Finite obstacle};
        \draw[<->,very thick,red!75!black] (1.15,0.25) -- (1.15,1.55)
            node[midway,right] {$z^\ast/R$};

        % Useful optical volume and adopted reference frame
        \draw[densely dotted,green!45!black,line width=0.8pt] (-4.30,-1.88) -- (-4.30,3.00);
        \draw[densely dotted,green!45!black,line width=0.8pt] ( 6.10,-1.88) -- ( 6.10,3.00);
        \coordinate (O) at (-6.25,0.25);
        \fill (O) circle (1.4pt);
        \draw[->,thick] (O) -- ++(0.95,0) node[right] {$x$};
        \draw[->,thick] (O) -- ++(0,0.95) node[above] {$z$};
        \draw[->,thick] (O) -- ++(-0.58,-0.42) node[below left] {$y$};
    \end{tikzpicture}%
    }
    \caption{Experimental setup, operating trajectory, Lighthouse coverage, and adopted reference frame.}
    \label{fig:experimental_setup}
\end{figure*}

During each run, flight data are streamed by radio to the ground station (position every 20~ms, the other channels every 50 or 100~ms), where synchronized records are stored at about 10~Hz. The logged quantities include position and translational velocity, attitude and angular rates, acceleration, battery voltage, collective-thrust command, and the individual motor pulse-width-modulation (PWM) commands. The complete record is retained for trajectory assessment and for reconstructing each physical episode in the DT. For GE identification, the physical and DT trajectories are aligned along the streamwise coordinate (Section~\ref{sec:information-phase}). The other logged channels serve trajectory assessment and do not enter the identification objective.

The Crazyflie is operated using its standard cascaded PID control architecture, as documented by Bitcraze~\cite{BitcrazeControllers}. During the crossing, the ground station sends a forward-velocity command scheduled on the measured $x$, a lateral velocity command that holds $y=0$, and the altitude setpoint, every 20~ms. An outer position loop converts the altitude setpoint into a vertical velocity reference. A velocity loop then regulates the translational motion, producing the desired roll and pitch angles together with the collective thrust command. These references are passed to the attitude controller, which determines the desired body angular rates. Finally, the inner angular-rate controller generates the roll, pitch, and yaw control contributions, which are combined with the collective thrust through the onboard motor mixer to determine the four individual motor commands. The DT reproduces this command mode, the firmware logic~\cite{BitcrazeFirmware}, and the loop rates (100~Hz for position and velocity, 500~Hz for attitude and rates), so that the DT and the vehicle run the same controller. The flight script aborts a mission when the position log is older than 0.35~s or the end of the mission is not reached in time, and it rejects reference corrections larger than $4.55R$.

\paragraph*{Digital Twin}
\label{sec:digital_twin}

The DT follows the model of Richardson and Pretorius~\cite{Richardson2025}, validated by its authors in near-hover flight, and is tailored to the Crazyflie~2.1+ of the campaign. The model is implemented as a differentiable physics-based simulator using the JAX library in Python and is used for both model identification and control optimization. The vehicle kinematics, rigid-body dynamics, propulsion system, actuator dynamics, and cascaded controller together form the nominal starting DT. The rigid-body, inertia, rotor-coefficient, actuator, and controller parameters are inherited from Richardson and Pretorius. The vehicle mass was measured, and the effective command-to-thrust gain $T_{\max}=0.1148$~N per motor was calibrated from the measured mass and the steady-hover collective command before the obstacle. No other parameter is reidentified. The finite-obstacle ground-effect correction is incorporated below as the unknown model component with associated $\boldsymbol{\theta}_p$ learned by RT.

The Crazyflie is modeled as a six-degree-of-freedom rigid body. Its state is described by position $\mathbf{x}=[x,y,z]^T$ and translational velocity $\mathbf{v}=[v_x,v_y,v_z]^T$ expressed in the inertial frame, the attitude quaternion $\mathbf{q}=[q_w,q_x,q_y,q_z]^T$, and the angular velocity $\boldsymbol{\omega}=[\omega_x,\omega_y,\omega_z]^T$ expressed in the body-fixed frame. The rotational speed of rotor $i$ is represented by $\Omega_i$.

The translational dynamics evaluated at the center of mass (CoM) are
%
\begin{equation}
\begin{cases}
{\dot{\mathbf{x}}=\mathbf{v}}\\
{m\dot{\mathbf{v}}=m\mathbf{g}+\mathbf{R}\mathbf{F}_{B}}
\end{cases}
\label{eq_vehicle_trasl}
\end{equation}
where $\mathbf{g}=[0,0,-g]^T$, $g$ is the gravitational acceleration, and $\mathbf{R}\in SO(3)$ transforms quantities from the body-fixed frame to the inertial frame. The total force in body coordinates is $\mathbf{F}_{B}=[0,0,T]^T$, where $T$ is the total thrust generated by the four rotors.
%

The rotational dynamics follow Euler's rigid-body equation,
%
\begin{equation}
\mathbf{I}\dot{\boldsymbol{\omega}}
+
\boldsymbol{\omega}
\times
(\mathbf{I}\boldsymbol{\omega})
= \boldsymbol{\tau}
\label{eq_vehicle_rot}
\end{equation}
where $\mathbf{I}$ is the inertia tensor and $\boldsymbol{\tau}=[\tau_x,\tau_y,\tau_z]^T$ is the total moment vector expressed in body coordinates.
%
The vehicle attitude is propagated through the unit quaternion according to
%
\begin{equation}
\dot{\mathbf q} = \frac{1}{2}
\mathbf{\Lambda}(\boldsymbol{\omega})\mathbf q
\label{eq_vehicle_attitude}
\end{equation}
where $\mathbf{\Lambda}(\boldsymbol{\omega})$ is the scalar-first quaternion kinematic matrix associated with the body angular velocity, using the convention adopted by Richardson and Pretorius~\cite{Richardson2025}.
%
The thrust and reaction torque generated by rotor $i$ are modeled as
%
\begin{equation}
\begin{cases}
{T_i=k_T\Omega_i^2}\\
{Q_i=k_Q\Omega_i^2}
\end{cases}
\label{eq_T_and_Q}
\end{equation}
where $T_i$ and $Q_i$ are the individual rotor thrust and torque and $k_T$ and $k_Q$ are the corresponding coefficients. The total thrust and moments follow from their signed combination.
%
For the adopted quadcopter and motor convention, the roll ($\tau_R$), pitch ($\tau_P$), and yaw ($\tau_Y$) moments are
%
\begin{equation}
\begin{cases}
{\tau_R=
\frac{d}{2}
\left(
-T_1-T_2+T_3+T_4
\right)}\\
{\tau_P=
\frac{d}{2}
\left(
-T_1+T_2+T_3-T_4
\right)}\\
{\tau_Y=
\sum_{i=1}^{4}s_iQ_i,}
\end{cases}
\label{eq_RP_moments}
\end{equation}
where $d$ is the distance between adjacent motors and $s_i\in\{-1,+1\}$ identifies the rotation direction of rotor $i$.
%

With the firmware battery compensation enabled, each motor command $c_i\in[0,65535]$ is a voltage-independent thrust request. The DT therefore maps it linearly to rotor thrust, $T_i=T_{\max}c_i/65535$, and sets the commanded rotor speed $\Omega_{i,\mathrm{cmd}}=(T_i/k_T)^{1/2}$. The realized speed follows this command through a first-order lag with time constant $\tau_{\mathrm{rpm}}=0.020$~s and a rate limit of $2\times10^5$~rpm/s.

Table~\ref{tab:dt_parameters} lists the DT parameters; the DT also uses the products of inertia of Richardson and Pretorius, up to $1.80\times10^{-6}$~kg~m$^2$. The dynamics are integrated with a fourth-order Runge--Kutta scheme at a 2~ms step, and aerodynamic drag is neglected.
\begin{table}[t]
\centering
\caption{DT parameters. The mass is measured, $R$ is the propeller radius, and $T_{\max}$ is calibrated from hover; the other values are those of Richardson and Pretorius~\cite{Richardson2025}.}
\label{tab:dt_parameters}
\begin{tabular}{lc}
\hline
Parameter & Value \\
\hline
$m$ & $0.03242$ kg\\
$R$ & $0.023$ m \\
$d$ & $0.065$ m \\
$I_{xx}$ & $16.57\times10^{-6}$ kg m$^2$ \\
$I_{yy}$ & $16.56\times10^{-6}$ kg m$^2$ \\
$I_{zz}$ & $29.26\times10^{-6}$ kg m$^2$ \\
$k_T$ & $2.88\times10^{-8}$ N s$^2$/rad$^2$ \\
$k_Q$ & $7.24\times10^{-10}$ N m s$^2$/rad$^2$ \\
$\tau_{\mathrm{rpm}}$ & $0.020$ s \\
$T_{\max}$ & $0.1148$ N \\
\hline
\end{tabular}
\end{table}
\paragraph*{Ground-effect correction}
\label{sec:ge-model}
The nominal DT does not account for the rotor--surface interaction generated as the vehicle crosses the finite obstacle at low clearance. This interaction is modeled with a low-order law inspired by the Cheeseman--Bennett relation for an isolated rotor~\cite{CheesemanBennett1955}. For a rotor $i$ of radius $R$ at a clearance $z_i$ from the obstacle, the in-ground-effect thrust $T_{i,\mathrm{IGE}}$ is obtained from the nominal out-of-ground-effect thrust $T_{i,\mathrm{OGE}}$ through the amplification factor $f_{\mathrm{GE}}$,
\begin{equation}
 T_{i,\mathrm{IGE}}
 =T_{i,\mathrm{OGE}}\,
 f_{\mathrm{GE}}(z_i;\boldsymbol{\theta}_p).
 \label{eq:ge-thrust-factor}
\end{equation}
The amplification factor is modeled as
\begin{equation}
 f_{\mathrm{GE}}(z_i;\boldsymbol{\theta}_p)=1+A\left(\frac{R}{z_i}\right)^p,
 \qquad \boldsymbol{\theta}_p=[A,\,p]^\top .
 \label{eq:fge}
\end{equation}
Here, $A$ controls the magnitude of the thrust augmentation, while $p$ determines how rapidly the interaction decays with increasing clearance. For a hovering rotor at constant power, Cheeseman and Bennett obtained $T_{\mathrm{IGE}}/T_{\mathrm{OGE}}=[1-(R/4z)^2]^{-1}$~\cite{CheesemanBennett1955}, whose first-order expansion gives $A=1/16$ and $p=2$. These values initialize the identification. The DT applies $f_{\mathrm{GE}}$ to the thrust at the actual rotor speed, not at constant power, so the identified $A$ and $p$ are effective coefficients of the vehicle in closed loop.

The model is static. It treats ground effect as a function of the instantaneous rotor clearance only. Kan et al.~\cite{Kan2019} found that, at low and medium speeds, ground effect depends much more on clearance than on forward speed. Up to the $0.5$~m/s flown here, the forward-flight term of the Cheeseman--Bennett relation also changes the augmentation by less than $2\%$. The model therefore neglects the speed dependence. It also omits the transient development and decay of the rotor--surface interaction at the obstacle edges~\cite{Curtiss1984}, and it is monotone in clearance, whereas a negative ground effect near $z/R\approx1.5$ has been reported for a closely spaced quadrotor~\cite{Lei2026}. The identified law is therefore an effective model over the tested clearances, not a general quadrotor law. These omissions preserve a compact model that can be identified and differentiated efficiently, while leaving a structural discrepancy that is addressed later through the reward-transfer surrogate.

Equation~\eqref{eq:fge} describes the clearance dependence above an extended surface. The obstacle considered here occupies only a finite streamwise interval, so the augmentation must also depend on the position of each rotor relative to its edges. The augmentation is therefore weighted by the fraction $\chi_i(x)\in[0,1]$ of rotor disk $i$ that overlaps the obstacle, in line with measurements on a partially covered rotor whose thrust correlates with the overlapping disk area~\cite{Yon2026}. The corrected rotor thrust becomes
\begin{equation}
 T_{i,\mathrm{corr}}=T_{i,\mathrm{OGE}}\left[1+\chi_i(x)\left(f_{\mathrm{GE}}(z_i;\boldsymbol{\theta}_p)-1\right)\right],
 \label{eq:corrected-thrust}
\end{equation}
where the overlap fraction is zero when the rotor lies outside the obstacle footprint and reaches unity when its disk is fully above it. Its continuous variation across the leading and trailing edges avoids discontinuous force changes and accounts for the four rotors entering and leaving the obstacle at different streamwise positions. It weights each rotor independently, whereas free-flight tests near a building-like edge found that full ground effect develops only after most of the vehicle has crossed the edge~\cite{Herz2026}. 

\subsection{Test Campaign Overview}
\label{sec:test-campaign}
Following Kan et al.~\cite{Kan2019}, the desired forward speed $v_x^\ast$ is normalized by the ideal induced velocity in hover,
\begin{equation}
 v_h=\sqrt{\frac{T_{h,r}}{2\rho_a\pi R^2}}
 =\sqrt{\frac{mg}{8\rho_a\pi R^2}},
 \label{eq:hover-induced-velocity}
\end{equation}
where $T_{h,r}=mg/4$ is the hover thrust of each rotor and $\rho_a$ is the air density. With $m$ and $R$ from Table~\ref{tab:dt_parameters} and $\rho_a=1.225$~kg/m$^3$, Eq.~\eqref{eq:hover-induced-velocity} gives $v_h\simeq4.42$~m/s. An operating condition is henceforth denoted by the ordered pair $(z^\ast/R,\,v_x^\ast/v_h)$.

The nominal operating condition is $(z^\ast/R,\,v_x^\ast/v_h)=(1.5,0.079)$. Unless stated otherwise, policy design, model identification, and reward evaluation refer to this condition. Table~\ref{tab:campaign} lists the policies flown and the number of flights. The campaign is organized into five consecutive blocks, introduced below together with the notation used in the results.

\begin{table}[t]
\centering
\caption{Policies and flights of the campaign. The record of each policy is the mean of its retained flights.}
\label{tab:campaign}
\begin{tabular}{lccc}
\hline
Stage and policies & Policies & Flights each & Flights \\
\hline
Baseline $\pi_{\mathrm{BL}}$ & 1 & 3 & 3 \\
Information $\pi_{\mathrm{INF}}$ & $10+2^{a}$ & 3--4 & 37 \\
Model-based $\pi_{\mathrm{MB}}$ & 1 & 3 & 3 \\
Interpolation $\pi_{\lambda,i}$ & 2 & 4 & 8 \\
Sobol $\pi_{\mathbf a,i}$ & 7 & 3--5 & 27 \\
EI $\pi_{\mathrm{EI},i}$ & 10 & 3--6 & 37 \\
At $(2.5,0.045)$ & 3 & 2--3 & 7 \\
At $(1,0.113)$ & 5 & 2--4 & 15 \\
\hline
Total & 41 & & 137 \\
\hline
\multicolumn{4}{l}{\footnotesize $^{a}$Archived records of an earlier campaign.}
\end{tabular}
\end{table}

\paragraph*{\textit{Baseline policy}}
The baseline policy $\pi_{\mathrm{BL}}$ is a constant reference at the nominal condition, tracked by the onboard cascaded PID controller. It is flown repeatedly to quantify run-to-run variability. Its mean response and RMS band provide the reference trajectory statistics, the first record in the replay buffer, and the baseline cost $J_{\mathrm{BL}}$ used in the normalized score of Eq.~\eqref{eq:normalized-score}.

\paragraph*{\textit{Information phase}}
\label{sec:information-phase}
The objective of the information phase is to design and perform trajectories that improve the identification of the unknown parameters $\boldsymbol{\theta}_p=[A,p]$ of the ground-effect model described by Eqs.~\eqref{eq:ge-thrust-factor} and~\eqref{eq:fge}.
Each episode is replayed in the DT with the deployed reference, starting at $x=1.70$~m from the steady out-of-ground-effect state at the commanded clearance and speed of the record. Factual replay identifies $A$ and $p$ from the accumulated tracking response, and the DT then designs the next experiment.

Starting from the baseline record, twelve information-oriented records $\pi_{\mathrm{INF}}^{(n)}$ are added to the cumulative dataset. Ten were designed in this campaign in four batches, the last of which kept a single policy. Six span the broad range $0.5\leq z/R\leq6$ of the reference, and four the low-clearance range $0.3\leq z/R\leq1.5$. The other two are low-clearance records of an earlier campaign at $z^\ast/R=1$ and $v_x^\ast/v_h=0.091$ and $0.102$, of which only the mean trajectories are available. The broad stage explores a wide set of clearance histories. It establishes the global scale and decay of the interaction without assuming that the most useful data come from the smallest admissible clearance, and it tests the evolving DT over the whole admissible range.

In the post-campaign evaluation of Section~\ref{sec:results-inf}, the broad-domain episodes after the first still increase the cumulative Fisher information, but each adds much less than the first did. This agrees with Eq.~\eqref{eq:fge}. At larger clearances, $f_{\mathrm{GE}}$ approaches unity and its sensitivity to both $A$ and $p$ decreases. After the sixth episode, the design domain of the reference is restricted to $0.3\leq z/R\leq1.5$, the low-clearance domain. All preceding episodes remain in the replay buffer, while the new candidates concentrate the available experimental effort where the identified law is steepest and the weak parameter direction can be excited more effectively. During the campaign, the model was updated after each batch. In the post-campaign identification of Section~\ref{sec:results-inf}, the records are appended one at a time, and each fit is warm-started from the preceding estimate.

The information phase begins with factual replay and adjoint identification. For physical episode $n$, the record provides the deployed reference and the measured clearance error $e_{z,n}^{\bullet}(x)=z(x)-z^\ast$. The DT replays the same episode with the current model parameters $\boldsymbol{\theta}_p$ and gives the corresponding error $e_{z,n}^{\circ}(x;\boldsymbol{\theta}_p)$. Both are evaluated at the $K_n$ samples of the measured record above the obstacle, $2.20\leq x\leq3.01$~m, and collected in the vectors
\begin{align}
 \mathbf e_{z,n}^{\bullet}
 &=\bigl[e_{z,n}^{\bullet}(x_{n,1}),\ldots,
 e_{z,n}^{\bullet}(x_{n,K_n})\bigr]^{\mathsf T},\\
 \mathbf e_{z,n}^{\circ}(\boldsymbol{\theta}_p)
 &=\bigl[e_{z,n}^{\circ}(x_{n,1};\boldsymbol{\theta}_p),\ldots,
 e_{z,n}^{\circ}(x_{n,K_n};\boldsymbol{\theta}_p)\bigr]^{\mathsf T}.
 \label{eq:episode-error-vectors}
\end{align}
With $N$ episodes in the replay buffer, the identification loss is
\begin{equation}
 J_p(\boldsymbol{\theta}_p)=
 \frac{1}{N}\sum_{n=1}^{N}
 \frac{1}{K_n}
 \left\|
 \frac{\mathbf e_{z,n}^{\circ}(\boldsymbol{\theta}_p)
       -\mathbf e_{z,n}^{\bullet}}{R}
 \right\|_2^2.
 \label{eq:identification-loss}
\end{equation}
The factor $1/K_n$ gives every episode equal weight irrespective of its number of samples or duration.

The loss is minimized by differentiating the complete time-discrete closed-loop rollout in reverse mode. This is the discrete-adjoint formulation of the identification problem: the loss sensitivity is propagated backward through the vehicle dynamics, actuator model, controller states, and ground-effect correction. At optimization step $t$, the resulting gradient is
\begin{equation}
 \mathbf{g}_t=
 \nabla_{\boldsymbol{\theta}_p}J_p
 \bigl(\boldsymbol{\theta}_{p,t}\bigr)
 =
 \begin{bmatrix}
 \dfrac{\partial J_p}{\partial A} &
 \dfrac{\partial J_p}{\partial p}
 \end{bmatrix}^{\mathsf T}.
 \label{eq:adjoint-gradient}
\end{equation}
It accounts for the direct influence of $A$ and $p$ on the aerodynamic forces and for the trajectory changes produced through the closed feedback loop. Both derivatives are obtained in one backward pass, without a separate perturbed simulation for each parameter.

Adam~\cite{Kingma2015} uses this gradient to update the parameters over successive optimization steps. It retains a first moment $\mathbf{m}_t$, which smooths the gradient direction, and a second moment $\mathbf{v}_t$, which tracks its component-wise squared magnitude:
\begin{align}
 \mathbf{m}_t &=
 \beta_1\mathbf{m}_{t-1}
 +(1-\beta_1)\mathbf{g}_t,\\
 \mathbf{v}_t &=
 \beta_2\mathbf{v}_{t-1}
 +(1-\beta_2)\mathbf{g}_t^{\odot 2}.
 \label{eq:adam-moments}
\end{align}
The two moments allow the step size to adapt separately to $A$ and $p$, whose gradients can have different scales. After the standard bias correction, the next parameter vector is
\begin{equation}
 \boldsymbol{\theta}_{p,t+1}=
 \Pi_{\Theta}\!\left[
 \boldsymbol{\theta}_{p,t}
 -\eta
 \frac{\widehat{\mathbf{m}}_t}
 {\sqrt{\widehat{\mathbf{v}}_t}+\epsilon}
 \right],
 \label{eq:adjoint-adam}
\end{equation}
where $\eta$ is the learning rate, $\epsilon$ prevents division by zero, and $\Pi_{\Theta}$ projects onto the admissible set $\Theta=\{0\leq p\leq6,\ A\geq-0.95\,(0.5)^{p}\}$, which keeps $f_{\mathrm{GE}}\geq0.05$ for $z_i\geq0.5R$. In the present campaign, $\eta=5\times10^{-3}$ and convergence is assessed over windows of 25 optimization steps. The relative-loss, coefficient-range, and projected-gradient tolerances are $2\times10^{-4}$, $10^{-3}$, and $5\times10^{-3}$, respectively. The optimization is considered converged only when all three conditions remain satisfied over a complete monitoring window, after at least 100 steps. When a new physical episode is added, the next fit is warm-started from the previously converged parameters, while the Adam moments are reset because the cumulative objective has changed. A cold start from $(A,p)=(0,0)$ on the complete dataset converges to $(0.120,1.970)$, against $(0.121,1.980)$ with warm starts.

The converged estimate $\widehat{\boldsymbol{\theta}}_p$ is then used to determine which experiment should be performed next. Local parameter uncertainty is approximated from the curvature of the identification loss~\cite{Ljung1999}. Let
\begin{equation}
 \mathbf H=\nabla^2J_p(\widehat{\boldsymbol{\theta}}_p)
 =\mathbf Q\operatorname{diag}(\lambda_1,\lambda_2)\mathbf Q^{\mathsf T}.
\end{equation}
To avoid a singular inverse along a locally flat direction, the eigenvalues are floored at $\lambda_H=10^{-6}$. The resulting local inverse-Hessian approximation is
\begin{equation}
 \begin{split}
 &\boldsymbol{\Sigma}_{\theta}\simeq
 \max\!\left[J_p(\widehat{\boldsymbol{\theta}}_p),10^{-12}\right]\mathbf Q\\
 &\quad\times\operatorname{diag}\!\left[
 \max(\lambda_1,\lambda_H)^{-1},
 \max(\lambda_2,\lambda_H)^{-1}
 \right]\mathbf Q^{\mathsf T}.
 \end{split}
 \label{eq:parameter-covariance}
\end{equation}
This matrix is used as a local curvature-based uncertainty approximation, not as a calibrated finite-sample confidence region. A flat direction of the loss produces a larger local uncertainty. For a model output $g(\boldsymbol{\theta}_p)$, such as $f_{\mathrm{GE}}$ at a prescribed clearance, the corresponding linearized uncertainty indicator is
\begin{equation}
 \sigma_g^2=
 \nabla_{\boldsymbol{\theta}_p}g^{\mathsf T}
 \boldsymbol{\Sigma}_{\theta}
 \nabla_{\boldsymbol{\theta}_p}g.
 \label{eq:uncertainty-propagation}
\end{equation}
The quantity $2\sigma_{f_{\mathrm{GE}}}/f_{\mathrm{GE}}$ reported later is the propagated relative uncertainty of the total thrust factor. It is not interpreted as an exact 95\% confidence interval, and its decrease at large clearance partly reflects that $f_{\mathrm{GE}}\rightarrow1$.

Candidate altitude-reference policies are then generated inside the admissible design space and simulated with the converged DT. Because $A$ and $p$ have different numerical scales, the sensitivity columns are normalized with the local parameter scales
\begin{equation}
 \mathbf D_{\theta}=\operatorname{diag}
 \left(\sqrt{[\boldsymbol{\Sigma}_{\theta}]_{AA}},
       \sqrt{[\boldsymbol{\Sigma}_{\theta}]_{pp}}\right),
 \qquad
 \overline{\mathbf S}_n=
 \frac{\partial\mathbf e_{z,n}^{\circ}}
 {\partial\boldsymbol{\theta}_p}\mathbf D_{\theta}.
\end{equation}
The eigenvalues can therefore be interpreted for perturbations of one local parameter standard deviation, rather than being dominated by the raw numerical scale of either coefficient. Under an independent, homoscedastic residual approximation, a common normalized residual variance is estimated as $\sigma_e^2=\max[J_p(\widehat{\boldsymbol{\theta}}_p),10^{-12}]$. The Fisher-information proxy accumulated from the $N$ completed experiments is
\begin{equation}
 \mathbf F_{1:N}=
 \sum_{n=1}^{N}
 \frac{\overline{\mathbf S}_n^{\mathsf T}\overline{\mathbf S}_n}
 {K_n\sigma_e^2},
 \label{eq:cumulative-fisher}
\end{equation}
where $K_n$ prevents a longer episode from appearing more informative solely because it contains more samples. Large eigenvalues of $\mathbf F_{1:N}$ correspond to locally scaled parameter combinations that the accumulated experiments constrain strongly; its smallest eigenvalue describes the least-observed joint parameter direction. Since spatial residuals can remain correlated, this matrix is a comparative metric rather than an exact statistical Fisher matrix. It is evaluated after the campaign, for every episode and for the repeated-baseline counterfactual, with the final DT, the final parameter scaling, and the same residual variance, so that the cumulative values are directly comparable.

Each virtual candidate policy $\pi$ produces a predicted Fisher contribution $\mathbf F(\pi)$. The next information-oriented policy is selected with the E-optimal criterion
\begin{equation}
 \pi_{\mathrm{INF}}^{(N+1)}=
 \arg\max_{\pi\in\Pi_{\mathrm{adm}}}
 \lambda_{\min}\!\left(
 \mathbf F_{1:N}+\mathbf F(\pi)
 \right),
 \label{eq:e-optimal-design}
\end{equation}
which is the idealized criterion. The campaign used a batch approximation of it. Each design iteration produced three policies from 68 candidates, 64 scrambled Sobol points and four structured shapes. The candidates were screened in the DT for reference rate ($|v_z|\leq0.5$~m/s), reference acceleration ($|a_z|\leq2$~m/s$^2$), completion of the rollout, and a lowest-rotor clearance of at least $0.3R$. The first policy maximizes $\lambda_{\min}$ of its own predicted Fisher matrix among the most novel half of the $30\%$ most informative candidates, where novelty is the distance of the predicted response from the measured trajectories. The second maximizes $\lambda_{\min}$ of the joint matrix with the first, and the third maximizes the novelty of the predicted response after the trailing edge, subject to minimum information levels. The predicted Fisher matrices are averaged over up to five sigma-point samples of $\boldsymbol{\theta}_p$. The accumulated information $\mathbf F_{1:N}$ of past experiments is not used in the design, where past trajectories enter through the novelty filter.

The identification is then repeated on the enlarged replay buffer, and the information phase ends with the criterion of Section~\ref{sec:rt-framework}.

\paragraph*{\textit{Model-based policy}}
\label{sec:mb-optimum}
After identification, $A$ and $p$ are frozen, and the DT is used to improve altitude tracking across the obstacle. The policy is a spatial reference correction, computed offline. During the flight, the ground station interpolates it at the measured streamwise position and sends the corrected altitude setpoint $z_{\pi}$ to the unchanged cascaded PID controller,
\begin{equation}
 z_{\pi}(x)=z_{\mathrm{BL}}(x)+\Delta z_{\pi}(x),
 \label{eq:policy}
\end{equation}
where $z_{\mathrm{BL}}$ is the nominal mission reference altitude which is constant along $x$. The correction modifies only the vertical setpoint, as a function of the measured streamwise position. The horizontal trajectory, attitude stabilization, motor mixing, and onboard PID gains remain unchanged. A spatial rather than temporal parameterization prevents small timing differences between DT and experiment from shifting the anticipatory command relative to the obstacle.

At the nominal operating condition, the reference correction is parameterized at $N_x=81$ nodes spaced by $25$~mm ($1.09R$) over $2.00\leq x\leq4.00$~m, linearly interpolated between them, and zero outside this interval. Its amplitude is bounded by $|\Delta z|\leq3R$, while the predicted lowest rotor must retain a clearance of at least $0.5R$ from the local surface.

The model-based policy $\pi_{\mathrm{MB}}$ is obtained offline from a finite-horizon nonlinear prediction of the complete closed-loop mission. At every optimizer step, the DT propagates the vehicle, propulsion, aerodynamic model, PID states, and reference correction over the entire prediction interval. Reverse-mode differentiation supplies the gradient with respect to all spatial reference values, and Adam updates their amplitudes subject to bound and endpoint projections. No optimization runs during the flight.

Performance is measured against the nominal reference $z_{\mathrm{BL}}=z^\ast$. The clearance $z_k$ is taken from the repetition-averaged trajectory at the $N_{\mathcal R}$ samples $x_k$ in the window $\Omega_{\mathcal R}=[2.20,3.70]$~m ($21.7\leq x/R\leq87.0$), with $N_{\mathcal R}\simeq42$ at the nominal condition. DT rewards are evaluated at the same samples by factual replay with the final coefficients. The tracking cost and reward are
\begin{equation}
 \begin{split}
 J_{\mathrm{trk}}(\pi)&=\frac{1}{N_{\mathcal R}}
 \sum_{x_k\in\Omega_{\mathcal R}}
 \left[\frac{z_k-z_{\mathrm{BL}}(x_k)}{R}\right]^2,\\
 \mathcal{R}(\pi)&=1-J_{\mathrm{trk}}(\pi).
 \end{split}
 \label{eq:reward}
\end{equation}
The reported policy is optimized with the single final identified DT; no model-uncertainty ensemble is used. The objective augments the tracking term with amplitude, curvature, and clearance penalties,
\begin{align}
 J_{\mathrm{MB}}=&\;J_{\mathrm{trk}}
 +\frac{0.10}{N_x}\left\|\frac{\Delta z_{\pi}}{R}\right\|_2^2 \notag\\
 &+\frac{2.0}{N_x-2}\left\|\mathbf{D}^2\frac{\Delta z_{\pi}}{R}\right\|_2^2
 +10^5\Phi_{\mathrm{clr}}(\pi),
 \label{eq:mb-objective}
\end{align}
where $J_{\mathrm{trk}}$ is evaluated over $2.20\leq x\leq4.00$~m, $\mathbf{D}^2$ is the second-difference operator on the nodes, and $\Phi_{\mathrm{clr}}$ is the mean squared violation of the $0.5R$ minimum clearance of the lowest rotor while the vehicle is above the obstacle. The amplitude term limits control authority, the curvature term suppresses spatially sharp commands, and the clearance term applies the $0.5R$ limit to each tilted rotor rather than only to the center of mass. Adam runs for 150 iterations with learning rate $10^{-3}$, starting from $\Delta z_{\pi}=0$. The final policy is replayed once more in the DT and rejected if the hard clearance check fails.

\paragraph*{\textit{Performance phase}}
\label{sec:performance-phase}
The performance search is conducted at the nominal operating condition with the identified parameters frozen. Two interpolation policies, denoted collectively by $\pi_{\lambda,i}$ with $\lambda\in\{1/3,2/3\}$, are first tested between $\pi_{\mathrm{BL}}$ and $\pi_{\mathrm{MB}}$ to establish the global relation between the DT and physical rewards. The search is then expanded around $\pi_{\mathrm{MB}}$ using three smooth spatial basis functions. A scrambled Sobol library~\cite{Owen1998} of $2048$ coefficient vectors is screened in the DT, and seven admissible policies $\pi_{\mathbf a,i}$, chosen to span different amplitude combinations, are flown to train the GP discrepancy. Ten policies $\pi_{\mathrm{EI},i}$ are then selected by EI and flown, and each physical reward is added to the surrogate dataset. The policy with the largest nominal-condition physical reward is denoted by $\pi_{\mathrm{EI}}^{\star}$.

A candidate policy is written as
\begin{equation}
 \Delta z_{\pi}(x)=\Delta z_{\mathrm{MB}}(x)+
 R\,\mathbf{b}^{\mathsf T}(x)\mathbf a,
 \label{eq:rt-policy}
\end{equation}
where $\mathbf a=[a_A,a_B,a_C]^\top\in[-1,1]^3$ and $\mathbf{b}(x)$ holds three basis functions. Each is a Gaussian centered at $x=2.50$, $2.95$, and $3.30$~m ($x/R\simeq35$, $54$, and $70$ in the shifted coordinate of the figures), with standard deviation $0.25$, $0.20$, and $0.25$~m. It is multiplied by $\sin^2[\pi(x-x_0)/L]$, with $x_0=L=2.00$~m, and normalized to unit peak, so that $a_j$ is a peak amplitude in rotor radii. The reference and rotor-clearance checks used for $\pi_{\mathrm{MB}}$ define the admissible set $\mathcal A_{\mathrm{adm}}$.

The affine component of the reward-transfer surrogate is fitted by ordinary least squares to the paired DT and physical rewards. The residual is modeled by a zero-mean Mat\'ern-$3/2$ GP. At the nominal operating condition, its input is $\mathbf a$, with a length scale of $0.30$ in each coordinate. The signal variance is set to the empirical variance of the affine residuals, with a lower bound of $0.05^2$, and no hyperparameter is optimized. The rewards are computed from the repetition-averaged trajectories and are treated as deterministic observations; a diagonal jitter of $10^{-8}$ is added only for numerical stability. Accordingly, the GP posterior standard deviation quantifies interpolation uncertainty within this surrogate and not run-to-run experimental dispersion. The resulting relation is
\begin{equation}
 \mathcal{R}_{\bullet}=
 \gamma_0+\gamma_1\mathcal{R}_{\circ}
 +\delta(\mathbf a,\boldsymbol{\chi})+\varepsilon,
 \label{eq:reward-transfer}
\end{equation}
where $\delta$ is the GP discrepancy and $\boldsymbol{\chi}$ contains optional operating conditions expressed by $(z^\ast/R,\,v_x^\ast/v_h)$. The posterior mean and variance of $\delta$ are updated after every performance-oriented episode. Expected Improvement is used with zero exploration margin, $\zeta=0$, and selects the next experiment,
\begin{align}
 Z(\mathbf a)&=
 \frac{\mu_{\bullet}(\mathbf a)-\mathcal{R}_{\bullet}^{+}}
 {\sigma_{\bullet}(\mathbf a)},\\
 \operatorname{EI}(\mathbf a)&=
 \left(\mu_{\bullet}-\mathcal{R}_{\bullet}^{+}\right)\Phi(Z)
 +\sigma_{\bullet}\phi(Z),\\
 \mathbf a_{n+1}&=
 \arg\max_{\mathbf a\in\mathcal A_{\mathrm{adm}}}
 \operatorname{EI}(\mathbf a),
 \label{eq:ei}
\end{align}
with the posterior variance floored at $10^{-8}$. EI is maximized over the Sobol library, six points at $\pm0.25$ on the axes, the coefficients of the policies already flown, and 128 Gaussian perturbations ($\sigma=0.10$) of the four best, all screened in the DT.

For transfer across operating conditions, the DT and physical performance are normalized by the corresponding condition-specific baseline cost,
\begin{equation}
 \mathcal S=1-\frac{J_{\pi}}{J_{\mathrm{BL}}}
   =1-\frac{1-\mathcal R_{\pi}}{1-\mathcal R_{\mathrm{BL}}},
 \label{eq:normalized-score}
\end{equation}
and $\boldsymbol{\chi}$ is augmented with the operating condition. To transfer a policy, each coefficient is scaled by the ratio $s_j$ of the root-mean-square amplitude of the condition-specific model-based correction to that of the nominal one, both weighted by $b_j$, $a_j^{\mathrm{tr}}=s_j a_j^{\star}$, where $\mathbf a^{\star}$ holds the coefficients of $\pi_{\mathrm{EI}}^{\star}$. A correction learned where the model-based correction is large is therefore reduced where the local model-based correction is small, rather than being transferred as a fixed altitude offset.

\paragraph*{\textit{Generalization}}
The final block evaluates whether the policy structure learned at the nominal condition remains useful when the operating condition changes. Two off-nominal conditions are tested, $(z^\ast/R,\,v_x^\ast/v_h)=(2.5,0.045)$ and $(1,0.113)$. At each, the baseline policy $\pi_{\mathrm{BL}}$ is flown, and a model-based policy $\pi_{\mathrm{MB}}$ is computed with the frozen DT. The best nominal RT correction is then transferred in the coordinates of the condition-specific model-based reference correction and denoted by $\pi_{\mathrm{EI}}^{\mathrm{tr}}$. Where an additional condition-specific EI search is performed, its best physical policy is denoted by $\pi_{\mathrm{EI}}^{\star}$. These tests separate direct transfer of the nominal correction from a local RT search at a new operating condition.


\section{Results and Discussion}
\label{sec:results}

\subsection{Baseline Response}
\label{sec:results-bl}

Fig.~\ref{fig:stage0} shows the baseline crossing at the nominal condition, flown with the onboard cascaded PID controller alone. The dashed line is the constant reference $\pi_{\mathrm{BL}}$. The controller acts only through feedback, without anticipating the aerodynamic interaction.

The solid line in Fig.~\ref{fig:stage0} reports the mean experimental trajectory $\mathrm{EXP}_{\mathrm{BL}}$ obtained from three repetitions, while the surrounding shaded region represents the corresponding $\pm\mathrm{RMS}$ run-to-run dispersion. After the repetitions are aligned at the leading-edge crossing and interpolated on the samples of the first one, this quantity is $\mathrm{RMS}(x)=[N_r^{-1}\sum_{j=1}^{N_r}(z_j(x)-\overline z(x))^2]^{1/2}$, with $N_r=3$. The band therefore describes repeatability about the measured mean and is not an uncertainty interval for that mean. Before the obstacle, the vehicle holds a nearly constant clearance close to the commanded value. Over the obstacle, it first rises to about $z/R=2.5$ and then settles around $z/R=2$. The additional thrust generated in ground effect produces the initial overshoot. The feedback controller reacts by reducing the motor speeds and commanded thrust, but this corrective action is insufficient to recover $z^*/R=1.5$ while the vehicle remains over the obstacle. The cascaded PID is tuned around hovering and contains no model that can anticipate the surface-induced thrust augmentation.

The opposite transient occurs at the trailing edge. Ground-effect augmentation disappears while the motor command is still reduced by the preceding feedback action, producing an undershoot down to about $z/R=0.4$. The vehicle then recovers the commanded clearance, to within $0.1R$, about $35R$ downstream of the trailing edge. Both excursions exceed the run-to-run dispersion, so they reflect a repeatable interaction rather than experimental scatter.

\begin{figure}[t]
\centering
\resultfigure{Results/stage_0_pid_baseline.pdf}
\caption{Baseline policy $\pi_{\mathrm{BL}}$ and mean experimental response $\mathrm{EXP}_{\mathrm{BL}}$ over three repetitions. The shaded band represents the corresponding $\pm\mathrm{RMS}$ dispersion.}
\label{fig:stage0}
\end{figure}

\subsection{Information Phase}
\label{sec:results-inf}

Each information-oriented reference is defined by its values at fixed streamwise nodes and interpolated between them. The design domain bounds this reference, not the trajectory, which departs from it under the closed-loop dynamics. The first six policies are designed in the \emph{broad domain}, $0.5\leq z/R\leq6$, and the following ones in the \emph{low-clearance domain}, $0.3\leq z/R\leq1.5$, where the ground-effect response is strongest and most sensitive to $A$ and $p$. Fig.~\ref{fig:informative-policies} shows representative references and trajectories. The broad-domain policies $\pi_{\mathrm{INF},1}$ and $\pi_{\mathrm{INF},6}$ request different clearance histories across the obstacle, whereas the low-clearance policy $\pi_{\mathrm{INF},9}$ keeps the crossing at smaller clearances. These trajectories are the ones used in factual replay and identification.

\begin{figure}[t]
\centering
\resultfigure{Results/stage_1_representative_informative_trajectories.pdf}
\caption{Representative information-oriented policies and physical trajectories in the broad and low-clearance domains.}
\label{fig:informative-policies}
\end{figure}

The estimates of the model parameters $\boldsymbol{\theta}_p=[A,p]^{\mathsf T}$ are reported in Fig.~\ref{fig:adjoint-convergence}(a). Both coefficients start from the Cheeseman--Bennett values and change to $A=0.104$ and $p=1.34$ after the first identification, which uses only the baseline record (gray region). Each subsequent point is obtained after appending one new episode to the replay buffer and repeating the cumulative identification. Broad-domain episodes lie in the olive region and low-clearance episodes in the green region.

The E-optimal quantity $\lambda_{\min}(\mathbf F)$ is reported in Fig.~\ref{fig:adjoint-convergence}(b). It is the cumulative, locally scaled Fisher-information proxy available along the weakest joint parameter direction and is compared for the designed sequence, denoted by \emph{Designed $\pi_{\mathrm{INF}}$}, and the counterfactual sequence \emph{Repeated $\pi_{\mathrm{BL}}$}. The latter counts the measured baseline episode once per cumulative episode, with its sampling and residual level, so that it measures repeated-design information under independent noise rather than new experiments. At the first information-oriented mission, corresponding to cumulative episode 2, $\lambda_{\min}(\mathbf F)$ reaches $0.676$ for \emph{Designed $\pi_{\mathrm{INF}}$} and $0.060$ for \emph{Repeated $\pi_{\mathrm{BL}}$}. Under the adopted local scaling, the design score of the designed episode is therefore more than ten times larger along the least-observed direction.

Over the next five broad-domain episodes, $\lambda_{\min}(\mathbf F)$ rises only from $0.676$ to $1.000$. This value is dimensionless and set by the scaling of Section~\ref{sec:information-phase}; it is not a normalization. The GE-parameter estimates also continue to move and reach $A=0.218$ and $p=2.792$ at the end of this stage. Their variation is smaller than after the initial identification, although they are not yet stable. The correlation $\rho_{A,p}$ of $A$ and $p$, computed from $\boldsymbol{\Sigma}_{\theta}$, drops from $0.974$ for the baseline record to $0.783$ after the first informative episode. It then rises to $0.883$ over the remaining broad-domain episodes, as further broad-domain data mainly reinforce the same coupled direction.

The broad domain nevertheless provides a useful first identification of the ground-effect law, especially at larger clearance. The corresponding function, denoted by $f_{\mathrm{GE}}\,(0.5<z/R<6)$, is compared with the Cheeseman--Bennett initialization and the final estimate in Fig.~\ref{fig:fge-quality}(a). The local linearized uncertainty indicator $2\sigma_{f_{\mathrm{GE}}}/f_{\mathrm{GE}}$ is reported in Fig.~\ref{fig:fge-quality}(b). The olive dash-dot curve predicts a stronger effect than Cheeseman--Bennett at small clearance, but there the indicator is still $45\%$ at $z/R=1$. Above about $z/R=2$, it is below $4\%$ because the broad trajectories provide coverage there and the total factor approaches unity.

The low-clearance stage produces trajectories such as $\pi_{\mathrm{INF},9}$ in Fig.~\ref{fig:informative-policies}, which spend more of the crossing in the region where $f_{\mathrm{GE}}$ is most sensitive to the two coefficients.

\begin{figure}[t]
\centering
\resultfigure{Results/stage_1_sequential_adjoint_convergence.pdf}
\caption{Sequential identification of $A$ and $p$ and cumulative E-optimal information obtained with the designed policies and the repeated-baseline counterfactual.}
\label{fig:adjoint-convergence}
\end{figure}

The effect of this decision is visible in both panels of Fig.~\ref{fig:adjoint-convergence}. The slope of $\lambda_{\min}(\mathbf F)$ increases again when the low-clearance episodes are introduced, and its cumulative value rises from $1.000$ to $3.242$, of which $0.733$ comes from the two archived records. Repeating the baseline record over the same 13 cumulative episodes would reach only $0.387$. At the same time, $A$ and $p$ shift toward smaller values and stabilize from approximately cumulative episode 11 onward. The final estimates are $A=0.121$ and $p=1.980$.

The correlation then decreases to $\rho_{A,p}=0.818$, which is still high. This residual coupling follows from the power law and from the limited range of clearance that can be flown. The sensitivities $\partial f_{\mathrm{GE}}/\partial A=(R/z)^p$ and $\partial f_{\mathrm{GE}}/\partial p=A(R/z)^p\ln(R/z)$ are nearly proportional over a narrow range of clearance, and the second vanishes at $z=R$ and changes sign across it. Coordinated variations of $A$ and $p$ therefore produce nearly identical $f_{\mathrm{GE}}(z/R)$ curves. Additional data can shrink the set of plausible coefficient pairs without removing this intrinsic elongated direction completely. The E-optimal criterion has therefore increased the information along the weakest observable direction, while the two physical parameters remain partially coupled.

Returning to Fig.~\ref{fig:fge-quality}, the final model is close to the broad-domain estimate above $z/R=2$ and differs from it below, with $f_{\mathrm{GE}}=1.121$ instead of $1.218$ at $z/R=1$. In this region, the indicator $2\sigma_{f_{\mathrm{GE}}}/f_{\mathrm{GE}}$ falls from $45\%$ to $15\%$ at $z/R=1$ and from $9.9\%$ to $4.3\%$ at $z/R=1.5$. These observations and the replay of Fig.~\ref{fig:baseline-dt} are the basis for ending the information phase (Section~\ref{sec:rt-framework}). Obtaining additional information would require trajectories even closer to the obstacle, reducing the safety margin and increasing the likelihood of contact or hardware damage. The accessible clearance is also bounded by the physical vertical envelope of the vehicle, which prevents meaningful exploration arbitrarily close to the surface.

\begin{figure}[t]
\centering
\resultfigure{Results/stage_1_identification_quality.pdf}
\caption{Identified ground-effect function and propagated uncertainty after the six broad-domain experiments and after combining the broad and low-clearance domains.}
\label{fig:fge-quality}
\end{figure}

The final exponent, $p=1.980$, is close to the Cheeseman--Bennett value $p=2$, whereas the amplitude, $A=0.121$, is almost twice $1/16$. A ground effect stronger than the classical prediction has also been measured on a tethered Crazyflie~2.0~\cite{Carter2021} and on a small multirotor~\cite{Sharf2014}. This supports a system-specific information phase when a low-order model is transferred from an isolated hovering rotor to a complete Crazyflie crossing a finite obstacle.

The relevance of this improvement to control is assessed through the factual replay reported in Fig.~\ref{fig:baseline-dt}. The baseline experiment is simulated using the Cheeseman--Bennett coefficients, the coefficients identified using only the broad domain, and the final coefficients identified by combining the broad and low-clearance domains. With the Cheeseman--Bennett coefficients, the DT predicts the rise over the obstacle and the loss of altitude downstream, but underestimates both. Both identified models predict a larger rise over the obstacle and a stronger post-obstacle descent, producing a closer match with the experiment.

The broad-domain and final replays of $\pi_{\mathrm{BL}}$ are nearly identical. The model obtained after the broad domain therefore already reproduces the dominant baseline response, and the low-clearance episodes refine $f_{\mathrm{GE}}$ below $z/R=2$ without changing this replay.

\begin{figure}[t]
\centering
\resultfigure{Results/stage_1_baseline_with_converged_dt.pdf}
\caption{Baseline experiment replayed with the classical initialization, the DT identified using only the broad domain, and the final DT identified using the broad and low-clearance domains.}
\label{fig:baseline-dt}
\end{figure}

The information phase thus replaces the generic Cheeseman--Bennett law with a law identified for the tested vehicle and clearance range.

\subsection{Model-Based Policy}
\label{sec:results-mb}

Fig.~\ref{fig:mb-optimum} compares the baseline policy with the model-based policy $\pi_{\mathrm{MB}}$, obtained by minimizing $J_{\mathrm{MB}}$ of Eq.~\eqref{eq:mb-objective}. Unlike the constant reference $\pi_{\mathrm{BL}}$, $\pi_{\mathrm{MB}}$ begins to reduce the commanded clearance before the leading edge and over the first part of the obstacle. This anticipates the thrust augmentation produced by ground effect. The reference is subsequently raised around the trailing edge, prompting the feedback controller to increase thrust before the aerodynamic contribution is removed. Farther downstream, the correction gradually returns to zero as the vehicle recovers the nominal clearance.

These anticipatory actions lower the measured trajectory over the obstacle, attenuate the overshoot, and shorten the downstream recovery. An undershoot nevertheless remains after the trailing edge, to $z/R\approx0.5$ against $0.4$ under $\pi_{\mathrm{BL}}$. The improvement therefore comes mainly from reducing the excursion over the obstacle and shortening the downstream tracking transient, rather than eliminating every local extremum.

With the reward of Eq.~\eqref{eq:reward}, the DT reward increases from $\mathcal R_{\circ}=0.676$ for $\pi_{\mathrm{BL}}$ to $0.968$ for $\pi_{\mathrm{MB}}$, and the physical reward from $\mathcal R_{\bullet}=0.544$ to $0.855$. Both rewards rank $\pi_{\mathrm{MB}}$ above $\pi_{\mathrm{BL}}$, so the DT is accurate enough to design a policy that improves physical tracking, and the gap between the two rewards measures the remaining transfer error.

The DT prediction $\mathrm{DT}_{\mathrm{MB}}$ in Fig.~\ref{fig:mb-optimum} shows the same gap. Over the obstacle, $\mathrm{DT}_{\mathrm{MB}}$ and $\mathrm{EXP}_{\mathrm{MB}}$ follow similar trends, although the virtual response still underestimates the measured overshoot. The largest difference occurs at the trailing edge. The DT predicts that raising the reference almost suppresses the downstream undershoot, whereas the vehicle still descends to $z/R\approx0.5$. The identification loss covers only the samples above the obstacle, so the DT is not fitted to this post-edge response. The DT therefore predicts the correct spatial structure and benefit of the compensation but overestimates its effectiveness in the transient regions where the static GE model is least representative.

\begin{figure}[t]
\centering
\resultfigure{Results/stage_2_mb_optimum.pdf}
\caption{Experimental baseline and model-based policies, together with the DT prediction of $\pi_{\mathrm{MB}}$.}
\label{fig:mb-optimum}
\end{figure}

\subsection{Performance Phase}
\label{sec:results-transfer}
\label{sec:results-exploration}
\label{sec:results-ei}

The performance phase learns the GP discrepancy between DT and physical rewards. Section~\ref{sec:results-mb} showed that the DT is optimistic around the trailing edge, where the static ground-effect law cannot represent the complete transient response. The reward-transfer surrogate accounts for the control-relevant part of this error without enlarging the model.

The initial reward-transfer surrogate dataset contains the two performance-oriented policies already available, $\pi_{\mathrm{BL}}$ and $\pi_{\mathrm{MB}}$, for which both DT and physical rewards have been evaluated. Fig.~\ref{fig:intermediate-policies} shows the two interpolation policies $\pi_{\lambda,i}$ of Section~\ref{sec:performance-phase}, which apply one and two thirds of the model-based correction and calibrate the affine part of the surrogate.

\begin{figure}[t]
\centering
\resultfigure{Results/stage_3_intermediate_policy_references.pdf}
\caption{Intermediate reference policies $\pi_{\lambda,i}$ between the baseline and model-based policies.}
\label{fig:intermediate-policies}
\end{figure}

The calibration is then refined near $\pi_{\mathrm{MB}}$. Three coefficients keep the number of flights needed to train the GP small; a richer correction could describe more shapes but would need more flights.

The basis centers A and B lie above the obstacle and C downstream of it (Fig.~\ref{fig:sobol}). Fig.~\ref{fig:sobol} shows representative samples of the Sobol library (faded blue) and the seven policies flown (dark-blue dashed). They combine different amplitudes and include changes that look unfavorable, such as raising the reference above the obstacle or lowering it downstream, as well as changes that extend the correction of $\pi_{\mathrm{MB}}$. This spread is required to teach the surrogate which local shape changes improve the experiment and which merely appear favorable in the DT.

\begin{figure}[t]
\centering
\resultfigure{Results/stage_4_sobol_and_tested_policy_references.pdf}
\caption{Representative unscreened samples of the Sobol library ($|a_j|\leq0.8$), basis locations, and the seven Sobol policies $\pi_{\mathbf a,i}$ flown.}
\label{fig:sobol}
\end{figure}

The reward pairs obtained after this local exploration are reported in Fig.~\ref{fig:exploration-transfer}. As expected, the two interpolated policies lie approximately along the trend joining $\pi_{\mathrm{BL}}$ and $\pi_{\mathrm{MB}}$. Most Sobol policies have DT rewards close to that of $\pi_{\mathrm{MB}}$, whereas their physical rewards range from $0.771$ to $0.932$, which reveals a policy-dependent GP discrepancy. These measurements calibrate the GP discrepancy near $\pi_{\mathrm{MB}}$. Once measurements exist on different sides of the model-based policy, the GP interpolation uncertainty can guide the next experiment through Expected Improvement.

\begin{figure}[t]
\centering
\resultfigure{Results/stage_4_reward_transfer_after_policy_exploration.pdf}
\caption{Reward-transfer surrogate after local policy exploration.}
\label{fig:exploration-transfer}
\end{figure}

The EI-selected policies are superposed on the calibrated reward-transfer surrogate in Fig.~\ref{fig:ei-transfer}. EI does not seek a reward increase at every iteration. Some policies are selected for their GP uncertainty, and their physical reward can remain below the best reward $\mathcal R_{\bullet}^{+}$ observed so far. The ninth of the ten EI-selected policies is $\pi_{\mathrm{EI}}^{\star}$. Relative to $\pi_{\mathrm{MB}}$, its DT reward increases only from $0.968$ to $0.978$, whereas the physical reward rises from $0.855$ to $0.952$. The DT alone ranks the two policies as nearly equivalent, and the physical rewards, through the GP discrepancy, reveal the difference. The fifth EI policy, whose coefficients differ from those of $\pi_{\mathrm{EI}}^{\star}$ by at most $0.024$, reached $0.947$, so the gain is not tied to a single policy. The coefficients of $\pi_{\mathrm{EI}}^{\star}$ lie at least $0.30$ from the faces of the box $[-1,1]^3$, which therefore does not bind at the best policy. The best Sobol policy reached $0.932$.

\begin{figure}[t]
\centering
\resultfigure{Results/stage_5_reward_transfer_after_EI.pdf}
\caption{Reward-transfer surrogate after the EI-selected physical evaluations; the star denotes the best policy $\pi_{\mathrm{EI}}^{\star}$.}
\label{fig:ei-transfer}
\end{figure}

The physical meaning of this reward increase is reported in Fig.~\ref{fig:ei-trajectory}, where the baseline, model-based, and best EI policies are compared with their mean experimental responses. Above the obstacle, $\pi_{\mathrm{EI}}^{\star}$ maintains a lower reference over a longer streamwise interval than $\pi_{\mathrm{MB}}$, which reduces the remaining positive excursion. Near the trailing edge, the command then rises higher, to about $z/R=3.1$ against $2.5$ for $\pi_{\mathrm{MB}}$. The downstream minimum rises from $z/R\approx0.5$ to $1.1$, and the recovery toward $z^{\ast}/R=1.5$ is smoother. This smaller post-obstacle excursion separates $\pi_{\mathrm{EI}}^{\star}$ from $\pi_{\mathrm{MB}}$; the two predictive solutions remain more similar over the obstacle itself.

\begin{figure}[t]
\centering
\resultfigure{Results/stage_5_baseline_MB_and_best_EI_trajectories.pdf}
\caption{Desired and mean experimental trajectories for $\pi_{\mathrm{BL}}$, $\pi_{\mathrm{MB}}$, and $\pi_{\mathrm{EI}}^{\star}$.}
\label{fig:ei-trajectory}
\end{figure}

The same comparison is reconstructed from video frames in Fig.~\ref{fig:real-comparison}. Whereas the curves in Fig.~\ref{fig:ei-trajectory} are averages of three repetitions of each case, each row of Fig.~\ref{fig:real-comparison} follows one representative flight. The single-run paths show the same ordering as the mean trajectories. The model-based and EI responses remain closer above the obstacle, while the best EI policy preserves the nominal clearance more effectively after the trailing edge.

\begin{figure*}[t]
\centering
\includegraphics[width=\textwidth]{Results/Real_comparison.jpg}
\caption{Video-frame comparison of representative nominal-condition flights: (a) baseline policy $\pi_{\mathrm{BL}}$, (b) model-based policy $\pi_{\mathrm{MB}}$, and (c) best Expected-Improvement policy $\pi_{\mathrm{EI}}^{\star}$. The overlaid curves trace the measured vehicle path.}
\label{fig:real-comparison}
\end{figure*}

\subsection{Generalization Across Operating Conditions}
\label{sec:results-generalization}

The final stage tests whether the correction learned at the nominal condition remains useful when both the desired clearance and the forward speed change.

Fig.~\ref{fig:generalization-trajectories} shows the lower-speed, larger-clearance crossing at $(z^\ast/R,\,v_x^\ast/v_h)=(2.5,0.045)$ and the higher-speed, smaller-clearance crossing at $(1,0.113)$. The nominal best correction, $\mathbf a^{\star}=[-0.66,0.70,0.20]^{\mathsf T}$, is then transferred with the scaling of Section~\ref{sec:performance-phase}, with $\mathbf s=[0.41,0.53,0.43]^{\mathsf T}$ at $(2.5,0.045)$ and $[1.85,1.51,1.60]^{\mathsf T}$ at $(1,0.113)$.

In both operating conditions, the condition-specific model-based policy markedly improves tracking relative to the corresponding baseline policy. At $(z^\ast/R,\,v_x^\ast/v_h)=(2.5,0.045)$, the transferred policy provides a further, albeit modest, improvement. At the higher-speed, small-clearance condition $(1,0.113)$, however, its physical reward is slightly lower than that of the model-based policy. The transferred correction was learned from discrepancies measured only at the nominal condition. At $(1,0.113)$, the vehicle spends $1.6$~s over the obstacle instead of $2.3$~s, the transferred coefficients leave the box $[-1,1]^3$, and the policies operate at clearances where $f_{\mathrm{GE}}$ is least constrained. The discrepancy can therefore differ from the nominal one.

\begin{figure}[t]
\centering
\resultfigure{Results/stage_6_generalization_trajectory_comparison.pdf}
\caption{Generalization at (a) $(z^\ast/R,\,v_x^\ast/v_h)=(2.5,0.045)$ and (b) $(1,0.113)$: PID, local model-based, transferred RT, and condition-specific EI policies.}
\label{fig:generalization-trajectories}
\end{figure}

Rewards at different conditions are compared through the scores of Eq.~\eqref{eq:normalized-score}, $\mathcal S_{\circ}$ for the DT and $\mathcal S_{\bullet}$ for the experiment. Each is the fractional reduction of the tracking cost relative to the baseline policy at the same condition, for which it is zero.

In the condition-aware surrogate, the GP discrepancy depends on both the policy and the operating condition~\cite{KrauseOng2011},
\begin{equation}
 \widehat{\mathcal S}_{\bullet}=
 \widehat{\mathcal S}_{\bullet,\mathrm{aff}}(\mathcal S_{\circ})+
 \widehat{\delta}\!\left(\mathbf a,z^\ast/R,v_x^\ast/v_h\right).
 \label{eq:generalized-surrogate}
\end{equation}
Here, $\widehat{\mathcal S}_{\bullet,\mathrm{aff}}$ is the affine fit of $\mathcal S_{\bullet}$ on $\mathcal S_{\circ}$. As implemented, $\widehat{\delta}$ takes as standardized inputs the coefficients $a_j$ divided by the local amplitude of the model-based correction, the residual of their projection on the basis, $v_x^\ast$, and $z^\ast/R$. The scores are means over the flights of each policy. Their standard deviation over the flights, divided by the baseline cost and at least $0.01$, weights the affine fit and enters the GP as noise. Additional experiments over the velocity--clearance domain would train this surrogate to propose a condition-specific correction instead of a single transferred one.

The normalized physical score is compared with its DT counterpart in Fig.~\ref{fig:generalization-reward}. Black markers correspond to $(z^\ast/R,\,v_x^\ast/v_h)=(1.5,0.079)$, orange markers to $(2.5,0.045)$, and purple markers to $(1,0.113)$. The transferred-policy diamond improves slightly on the model-based-policy square at the low-speed, larger-clearance condition. At the higher-speed, small-clearance condition, the ordering is reversed and the transferred policy performs slightly worse than $\pi_{\mathrm{MB}}$. A correction calibrated at one operating condition therefore cannot be assumed to retain the same ranking elsewhere, which motivates extending the surrogate across the operating domain rather than relying on direct policy transfer alone.

\begin{figure}[t]
\centering
\resultfigure{Results/stage_6_PID_normalized_reward_generalization.pdf}
\caption{Normalized reward-transfer surrogate across the nominal and two additional operating conditions.}
\label{fig:generalization-reward}
\end{figure}

The higher-speed, small-clearance condition is therefore used to demonstrate how the normalized reward-transfer surrogate can be extended locally. Two EI rounds are run at $(z^\ast/R,\,v_x^\ast/v_h)=(1,0.113)$; each proposes three distinct candidates, of which one is flown. The first explores an uncertain policy direction and performs poorly, yielding approximately $\mathcal S_{\bullet}=0.57$, as indicated by the small purple marker in Fig.~\ref{fig:generalization-reward}. Once this observation updates the condition-aware GP discrepancy, the second evaluation reaches a new condition-specific best of approximately $\mathcal S_{\bullet}=0.83$, marked by the purple star, and exceeds the corresponding model-based policy. Its policy and physical trajectory are included in Fig.~\ref{fig:generalization-trajectories}(b). Unlike the nominal transferred correction, this policy remains closer to $\pi_{\mathrm{MB}}$, suggesting that the model-based policy was already near a favorable local optimum at the smaller clearance and higher velocity.

With only two EI iterations, this result is a proof of concept rather than a converged surrogate. At a new operating condition, the frozen DT first supplies a model-based policy that improves on the baseline policy. A few physical episodes can then reveal how the GP discrepancy has changed and redirect the search toward a better condition-specific policy.

\section{Conclusions and Outlook}
\label{sec:conclusion}

This work demonstrated Reinforcement Twinning for the trajectory control of a micro aerial vehicle crossing a finite obstacle in ground effect. In the information phase, the designed experiments raised the smallest eigenvalue of the cumulative Fisher-information proxy to $3.242$ over the baseline and twelve information-oriented episodes, against $0.387$ for the baseline record repeated over the same episodes. The low-clearance records, two of them archived, gave most of this increase. The parameters remain correlated ($\rho_{A,p}=0.818$), but the uncertainty indicator of $f_{\mathrm{GE}}$ fell to $15\%$ at $z/R=1$ and $4.3\%$ at $z/R=1.5$.

The DT produced a model-based policy that raised the physical reward from $0.544$ to $0.855$, kept the onboard feedback controller, and passed the DT clearance checks. Interpolation and Sobol policies then calibrated the reward-transfer surrogate, and after nine EI-selected policies the search found a policy with physical reward $0.952$, mainly by reducing the post-obstacle undershoot that the static model does not predict.

The off-nominal experiments clarified both the potential and the present limits of generalization. Directly transferring the nominal correction produced a modest gain at one operating condition but slightly underperformed the corresponding model-based policy at the other. After two condition-specific EI evaluations, the normalized reward-transfer surrogate identified an improved local policy. This result supports the data efficiency of local refinement, but it does not establish a surrogate that is converged over a continuous velocity-clearance domain.

The results are limited to a static two-parameter ground-effect law, to the clearances flown (the lowest measured mean clearance over the obstacle was $0.68R$), and to two off-nominal conditions. Admissibility was checked only in the DT, which underpredicts the trailing-edge undershoot, so the DT screening reduces the risk of a policy without bounding it. Future work should cover a denser velocity--clearance domain, include repeated-measurement uncertainty in the surrogate, and test unsteady aerodynamic models, moving landing surfaces, and onboard or real-time policy adaptation.

\section*{Acknowledgment}
The authors gratefully acknowledge the technical support of the von Karman Institute for Fluid Dynamics.

\bibliographystyle{IEEEtran}
\bibliography{references}

\end{document}



